\documentclass[11pt,letterpaper]{article}
\usepackage[T1]{fontenc}
\usepackage[utf8]{inputenc}
\usepackage{lmodern,microtype}
\usepackage[margin=0.87in,headheight=15pt]{geometry}
\usepackage{amsmath,amssymb,amsthm,mathtools}
\usepackage{booktabs,longtable,array,tabularx,enumitem}
\usepackage{xcolor,fancyhdr,listings,xurl,needspace}
\usepackage[colorlinks=true,linkcolor=blue!35!black,citecolor=blue!35!black,urlcolor=blue!45!black]{hyperref}
\usepackage{bookmark}
\definecolor{ink}{RGB}{24,43,62}
\definecolor{pale}{RGB}{245,247,249}
\setlength{\emergencystretch}{3em}
\setlength{\parskip}{0.3em}
\setlist{itemsep=3pt,topsep=5pt,leftmargin=1.6em}
\newtheorem{theorem}{Theorem}[section]
\newtheorem{lemma}[theorem]{Lemma}
\newtheorem{proposition}[theorem]{Proposition}
\newtheorem{corollary}[theorem]{Corollary}
\theoremstyle{definition}\newtheorem{definition}[theorem]{Definition}
\theoremstyle{remark}\newtheorem{remark}[theorem]{Remark}
\newcommand{\ipra}{\mathsf{iPRA}}
\newcommand{\D}{\mathsf D}
\newcommand{\Q}{\mathsf Q}
\newcommand{\Eq}{\mathsf{Eq}}
\newcommand{\Tr}{\mathsf{Tr}}
\newcommand{\FV}{\operatorname{FV}}
\newcommand{\llbracket}{\mathopen{[\![}}
\newcommand{\rrbracket}{\mathclose{]\!]}}
\newcommand{\file}[1]{\expandafter\nolinkurl\expandafter{\detokenize{#1}}}
\newcommand{\Acal}{\mathcal A}
\newcommand{\Ccal}{\mathcal C}
\lstset{basicstyle=\ttfamily\footnotesize,columns=fullflexible,keepspaces=true,
 breaklines=true,showstringspaces=false,frame=single,rulecolor=\color{black!15},
 backgroundcolor=\color{pale},aboveskip=8pt,belowskip=8pt}
\pagestyle{fancy}\fancyhf{}
\fancyhead[L]{\small iPRA: smaller checking and coverage without injectivity}
\fancyhead[R]{\small 10 November 2026}
\fancyfoot[C]{\thepage}
\hypersetup{pdftitle={Smaller iPRA Checking and Coverage Without Injectivity},pdfauthor={Research prepared for Vladimir Reshetnikov with ChatGPT}}
\title{\color{ink}\bfseries Smaller iPRA Checking\\and Coverage Without Injectivity\\[9pt]
\large A reproducible comparison and two arithmetic lemmas}
\author{Research prepared for Vladimir Reshetnikov with ChatGPT}
\date{10 November 2026}
\begin{document}
\maketitle
\begin{abstract}
The current ProveIt drafts are audited at a fixed commit, including nested
archives and complete ancillary checkers. A composition of existing scope,
arity and neutral-major techniques gives a complete initialized-iterator
iPRA checker of 529 normalized Python bytes, below the supplied 542-byte
incumbent; an explicit-list variant uses 537 bytes. This improvement does
not beat the independent token and AST minima, and is not a new arithmetic
foundation. The complete native axioms and rules, derived intuitionistic
logic, full-domain bi-interpretation and both internal composites are
provided with attribution.

The new mathematical contribution concerns the scalar self-seeded precursor
with $F_h(0)=S(h(0))$ and $F_h(Sn)=h(F_h(n))$. Actual Boolean equality and
separation exclude every fixed finite successor cycle even without
induction. With literal-equation induction, a guarded capped clock gives
a finite-case right predecessor and proves numerical coverage without
successor injectivity. The resulting identity is $S(P(Sn))=Sn$; the
left-inverse identity $P(Sn)=n$ is not proved. Accordingly, removing
injectivity from the known complete self-seeded iPRA theory remains open.
The requested significant new simplest axiomatization has not been achieved.
\end{abstract}
\begin{center}
\fcolorbox{ink}{pale}{\begin{minipage}{0.92\linewidth}
\textbf{Established:} a smaller complete checker for an inherited theory,
new coverage and finite-aperiodicity lemmas without injectivity, and a
reproducible source and certificate audit.\par
\textbf{Still open:} a significant new axiomatization below the full
complexity frontier, successor injectivity for the weakened self-seeded
precursor, and absolute minimality.
\end{minipage}}
\end{center}
\clearpage
\setcounter{tocdepth}{1}
\hypersetup{bookmarksdepth=2}
\begingroup\small\setlength{\parskip}{0pt}\tableofcontents\endgroup
\clearpage

\section{Results and their exact scope}
\label{sec:results}

The requested significant new axiomatization below the supplied complexity
frontier has not been established. This report gives concrete implementation
and arithmetic results, and keeps their claims separate.

First, a complete checker for the existing initialized scalar-iterator
presentation occupies 529 normalized Python bytes, compared with the current
542-byte incumbent. An explicit-list version occupies 537 bytes. Both include
the parser interface, scope checking, arithmetic, logic, and verdict. The
reduction is an implementation composition of techniques already in the
drafts, not a new arithmetic foundation. Neither program improves the independent
261-token and 291-AST-node minima, which are attained by different programs.

Second, a new capped-clock construction proves numerical coverage in the
scalar self-seeded precursor without successor injectivity. It defines a
finite-case graph $P$ and proves
\[
 n=0\ \lor\ n=S(Pn),\qquad S(P(Sn))=Sn.
\]
It does not prove $P(Sn)=n$. A separate argument rules out every fixed finite
successor cycle without any numerical induction. The complete self-seeded
iPRA presentation in the drafts still needs its injectivity axiom; these
theorems do not yet justify deleting it.

An initially promising syntax-directed paired binder turned out to be an
existing proposal in a nested archive. The earlier all-starts development
already gives deterministic introduction, the progress/base elimination
interface, and a fusion theorem. It is not counted as a new result here.
The independently reworked prototype was also larger than the incumbent.
This negative novelty check is recorded so that the same idea is not
mistaken for a future breakthrough.\cite{allstarts}

\subsection{The frozen comparison}

All 20 direct ZIPs in \file{docs/drafts} were retrieved at commit
\file{b9c4d74a0018bdafc5bc06e755a0145f4ed2b8d6}. Their exact byte lengths
and Git blob identities were verified. The recursive scan contains 35 distinct
archive payloads, 2,744 Python occurrences, and 484 distinct Python source
contents. All distinct source bytes are included for offline remeasurement;
the measurement command parses them but does not execute them. All 41
source files at or below any incumbent coordinate minimum were read in full:
three are complete checkers and 38 are supporting utilities.

The coordinate minima are achieved upper bounds on an unknown optimum.
They are not proved lower bounds on every possible presentation. A universal
lower-bound claim would additionally need a fixed implementation language,
runtime boundary, source convention, and a mathematical minimality argument.
No such claim is made here.

\begin{center}\small
\begin{tabular}{@{}lrrrr@{}}
\toprule
Complete program & Bytes & Tokens & AST & Gzip 9\\\midrule
Current byte incumbent & 542 & 263 & 309 & 343\\
Current explicit-list incumbent & 550 & 263 & 307 & 349\\
Current token incumbent & 578 & 261 & 295 & 346\\
Current AST incumbent & 642 & 263 & 291 & 353\\\midrule
New compact composition & \textbf{529} & 266 & 310 & 333\\
New explicit-list composition & \textbf{537} & 266 & 308 & 337\\
\bottomrule
\end{tabular}
\end{center}

Every row is one whole module. The compact composition improves the observed
byte minimum by $13/542\approx2.40\%$. The explicit-list version is five bytes
below that same minimum and thirteen below the matched explicit-list incumbent.
The gzip column compresses raw source at level nine, with \texttt{mtime=0}
and the wrapper included; it is informational only. Counts use CPython
3.12.14 and zlib 1.3.2. Normalization removes genuine docstrings (retaining an empty-suite
\texttt{pass} where needed), applies
\texttt{ast.unparse}, and appends a newline. Significant tokens omit whitespace,
comments, indentation and end markers. AST counts charge the whole normalized
module. Different runtimes may change normalization or compression.

\subsection{Proven, measured, inherited, and unresolved}

The initialized-iterator axiomatization, its full-domain iPRA bridge, the
arity codec, certified-zero repair, lexical scopes and neutral-major technique
are inherited and explicitly attributed. The smaller source composition is
new in this deliverable, but its ingredients are not claimed as novel.
The capped-clock right predecessor and fixed finite aperiodicity argument
are new results relative to the inspected material.
No literature-wide priority claim is made.

The metatheorems are written mathematical proofs, checked separately by other
assistant instances during this investigation. They are not Lean/Rocq
formalizations or external human peer review. Finite tests check executable
interfaces and arithmetic transcriptions; they do not establish the general
model-theoretic claims. No claim of a significant foundational breakthrough
or absolute simplestness is warranted by these results.

\section{The measured theory in a compact ordinary form}
\label{sec:ordinary-core}

For orientation, the following ordinary first-order theory $T_I$ is
bi-interpretable with full first-order intuitionistic PRA. Its numerical
terms are $x,0,St,I(n,a;v.h)$. Only $h$ binds the state variable $v$;
passive parameters are fixed. For each fixed earlier-stage update $h$, the
first-order realization has a function symbol $I_h(n,a,\vec p)$.
There is no numerical code ranging over a uniformly total evaluator.

Its nonlogical computation schemas are
\[
 I(0,a;v.h)=a,\qquad I(Sn,a;v.h)=h[I(n,a;v.h)/v].
\]
It has numerical coverage $x=0\lor\exists y\,x=Sy$, separation
$(0=S0)\to\bot$, and contextual induction on literal equations:
\[
 \frac{\Gamma\vdash E(0)\qquad\Gamma,E(z)\vdash E(Sz)}
      {\Gamma\vdash E(t)},
 \qquad E\text{ a literal numerical equation}.
\]
The parameter $z$ is fresh for $\Gamma$ and the passive parameters.
Intuitionistic first-order equality logic is part of this ordinary
presentation. The smaller native presentation, whose complete rules are
given in Section~\ref{sec:choice-introduction}, derives this logical interface,
coverage, separation and equation induction. These are not uncharged
primitive rules in the measured checker.

The native paired binder has the ordinary meaning
\[
 \mathsf D_{x,y}B\quad\mapsto\quad
 \exists x\bigl(B(x,0)\to\forall y B(x,y)\bigr).
\]
Equations and this binder are the only object-formula constructors. The
native primitives are computation, reusable hypothesis/closure use, paired
elimination, numerical transport, and case/progress paired introductions,
with progress permitted only at a literal equation. Their exact eigenvariable
conditions, derived logical macros, and formula composites appear below.

\paragraph{Why this is full iPRA.}
The dependency order matters. Equation induction and numerical cases first
construct a scalar predecessor, without assuming pairing, actual equality
decidability, or successor injectivity. Predecessor then gives truncated
subtraction and a characteristic of actual equality. Every fixed
quantifier-free formula becomes a characteristic equation, so all its
induction instances follow. Pairing and projections are then constructed
and every fixed primitive recursion is compiled into scalar iteration.
Internal uniqueness of the recurrences proves both numerical composites.
The domain remains the entire original numerical domain, and the definable
isomorphism in both directions is $x=y$. The attributed proofs are reproduced
in the appendices; agreement only on standard numerals would not suffice.

\section{What the 529-byte checker changes}
\label{sec:composition}

\subsection{The entire compact module}

\lstinputlisting[language=Python]{code/checker_compact.py}

\file{code/checker_explicit.py} uses \texttt{isinstance} to recognize list
nodes. This costs eight normalized bytes and is the recommended compact
source for direct inspection.
\file{code/readable_checker.py} is a separately executable expansion with
named operations and comments, not an omitted trusted dependency. No checker
imports a producer, numerical library, another kernel, or an oracle.

\subsection{Three compositions, with matched controls}

The current byte incumbent uses a global fresh-name floor. The new kernel
passes that floor lexically into each premise; completed major proofs cannot
change their caller's scope. This local discipline already exists in the
uniform-zero draft. The initial hypothesis is the empty-array truth, and
every packet checks its major, with \texttt{[]} serving as the immediately
proved neutral major. That protocol already exists in the monus/neutral-major
draft. Finally, each successful non-introduction branch constructs a nonempty
list, so Boolean disjunction can select the same hypothesis as an explicit
conditional chain. That factoring is also inherited.

\begin{center}\small
\begin{tabular}{@{}lrrrr@{}}
\toprule
Matched module & Bytes & Tokens & AST & Gzip 9\\\midrule
Guarded major, lexical floor & 545 & 267 & 311 & 336\\
Neutral major, global floor & 542 & 263 & 309 & 343\\
Neutral major, lexical, conditional dispatch & 535 & 266 & 310 & 335\\
Neutral major, lexical, factored dispatch & 529 & 266 & 310 & 333\\
\bottomrule
\end{tabular}
\end{center}

The controls isolate the change; they do not invent a new arithmetic basis
from incidental syntax savings. In particular, lexicalization followed by
factoring saves seven then six bytes against the 542-byte incumbent, while
adding three tokens and one AST node overall.

\subsection{The complete canonical packet interface}

The JSON input is an array. Its last entry is a closed goal; earlier
four-field packets $[m,s,u,t]$ are consumed backwards. A successful run
exits with status zero and prints exactly \texttt{ACCEPT}. An exception
rejects. An unused prefix has no effect. No Python code in a certificate
is executed.

Every non-hypothesis command first proves its major $m$ in the same context
and at the same lexical floor. A substantive major can transport equality,
specialize a closure, or eliminate a paired formula. A neutral major is
\texttt{[]}. With a neutral major and a nonempty initializer $u$, the
packet supplies a computation, with $s$ its predecessor discriminator and
$t$ its update. With $m,u$ both false-valued, it introduces a unary paired
goal with witness $t$. The marker $u=0$ selects equation progress; $u=[]$
selects the zero case. The unused $s$ field does not select induction.

The sorted representation uses unary arrays for successor or paired
binding, binary arrays for equations, and ternary arrays for scalar
iteration or auxiliary closure. A canonical closure is $[B,B,B]$, so
tripling a unary major exports precisely its closure. Other nonempty
array majors export a high-arity truth alias. At numerical positions,
empty and surplus arrays denote zero; at formula positions they denote
the derived truth $0=0$. These are explicit constant interpretations,
not hidden axioms about new symbols.

The common opening depth increment is $5-r$ at arity $r$. Numerical root
depth is three; one protected numerical slot has width two. Negative
admitted labels are rigid parameters. Under an iterator the counter and
initializer receive a dummy zero slot, while the update receives the
state slot. Under successor there are two dummy zero slots. The precise
decoder and opening proof are reproduced in
Appendix~\ref{sec:arity-checker}. This representation is deliberately
charged and explained; it is not ordinary surface arithmetic notation.

\subsection{Why all computations are still available}

Literal-zero initializers and literal-zero successor predecessors do not
select the compact direct-computation branch. Put
\[
                  Z_*:=I(S(S0),S0;v.0).
\]
Its permitted successor computation gives $Z_*=0$. In an excluded
computation instance, replace all offending initializer/predecessor positions
simultaneously by a fresh placeholder $w$, giving an equation template
$E(w)$. The equation $E(Z_*)$ is permitted. Transport along $Z_*=0$
then gives $E(0)$. This is the inherited three-node certified-zero repair.
It does not need to decide whether an arbitrary term denotes zero.

\subsection{Soundness of the composition}

Two details require more than a superficial source comparison. First,
the matcher has a Boolean-false alternative when $s$ is nonempty.
It would be unsafe if a numeric zero could be a proved major. The
neutral-major invariant excludes this: from list-valued assumptions,
no finite successful invocation proves a numeric scalar. A minimal-trace
argument proves the invariant for every starting scope floor. The initial
assumption is a list, and every surviving continuation adds a list. Thus
a false-valued major reaching the matcher is the empty list, never zero.
The complete inherited argument is in Appendix~\ref{sec:neutral-major-audit}.
Its quantification over starting floors makes it valid for lexical floors.

Second, generated computation hypotheses can contain freshly admitted
parameters. They need only have a provable simultaneous zero instance.
If $\vec b$ is fresh for $\Gamma,A$, from
$\vdash H(\vec0)$ and $\Gamma,H(\vec b)\vdash A$ one substitutes zero
for the new names and cuts in the computation proof. One must not claim
that every generated $H(\vec b)$ is uniformly true. Fractional numeric
aliases that pass the raw checker are included among the finite fresh
parameter family; canonical proofs use integer coordinates.

The remaining composition steps are direct. Premise contexts are fresh
lists, so they cannot mutate a caller. A nonempty list major gives
nonempty $m*3$; a computation produces a nonempty binary equation.
Therefore the factored dispatch has exactly the conditional dispatch's
value, and only introduction performs the assignment that opens the goal.
In introduction $u$ is false-valued and cannot equal odd arity one, so
compound progress is excluded. Canonical proof translation changes absent
majors to the empty list and moves the old choice marker into $u$; it
otherwise preserves the inherited native derivation.

Consequently every successful raw run decodes to an inherited derivation,
and every canonical inherited proof has a finite translated certificate
under adequate computational resources. Both iPRA interpretation
composites and equiconsistency transfer unchanged. Real stack and memory
limits can reject large valid certificates. The strict-list variants do
not use the generic malformed-string traversal that qualifies some
historical token/AST record holders.


% Standalone includable research fragment. Requires amsmath, amssymb,
% amsthm and theorem/lemma/corollary environments in the parent document.
% This is a new result in the present audit, not an assertion of priority
% outside the supplied repository corpus.
\section{Finite aperiodicity without injectivity or induction}
\label{sec:finite-aperiodicity}

The scalar self-seeded draft charges successor injectivity as an
axiom.\cite{selfseed}
There is a useful reduction of its unresolved arithmetic problem: every
\emph{fixed finite} successor cycle is already excluded by its computation
schemas and Boolean equality, without using either injectivity or any
induction rule. This does not exclude a merger of two distinct predecessor
chains. The distinction is essential.

\subsection{Exact assumptions}

Work in intuitionistic first-order logic with actual equality and falsity.
The numerical language has $0,S,\epsilon$, and, for each fixed earlier-stage
term $h(v,\vec p)$, a scalar iteration symbol $F_h(n,\vec p)$ satisfying
\begin{align}
 F_h(0,\vec p)&=S(h(0,\vec p)),\label{eq:ac-F0}\\
 F_h(Sn,\vec p)&=h(F_h(n,\vec p),\vec p).
 \label{eq:ac-FS}
\end{align}
The other assumptions used in this section are precisely
\begin{gather}
 0=S0\longrightarrow\bot,\label{eq:ac-sep}\\
 \epsilon(x,y)=0\ \lor\ \epsilon(x,y)=S0,
 \label{eq:ac-range}\\
 \epsilon(x,y)=S0\quad\longleftrightarrow\quad x=y.
 \label{eq:ac-reflection}
\end{gather}
The ordinary logical rules are specified in the next paragraph; they are
part of this precursor, not omitted from its cost. Write $1=S0$ and
$\bar j=S^j0$ for an externally fixed nonnegative integer
$j$. There is no quantification over function descriptions. Each displayed
$h$ is a finite term, and its $F_h$ symbol is introduced at a later finite
stage. Numerical coverage, successor injectivity, predecessor, arbitrary
initialization, and induction are not assumptions of this section.

\paragraph{Ordinary logical rule inventory.}
In the precursor, formulas are generated from numerical equations and
$\bot$ by $\land,\lor,\to,\forall,\exists$. Contexts are finite reusable
assumption lists. The complete logical inventory is hypothesis membership;
conjunction introduction and its two projections; the two disjunction
injections and disjunction elimination with a common conclusion from its
two discharged branch assumptions; implication introduction by discharging
its antecedent and modus ponens; universal introduction at a parameter
fresh for the retained context and specialization at a term; existential
introduction with a term witness and elimination at a parameter fresh for
the retained context and conclusion; and inference of any formula from
$\bot$. Equality has reflexivity and capture-free replacement of an equal
numerical term in an arbitrary formula. Weakening, exchange, contraction,
renaming and cut are admissible structural rules. These ordinary rules
have no excluded-middle or double-negation-elimination axiom. Quantifier
substitutions avoid capture. The displayed computation, separation,
range and reflection formulas are the only nonlogical axioms used here;
the next section adds exactly contextual literal-equation induction.

\begin{lemma}[Zero is not a successor]
\label{lem:ac-zero}
For every $a$, the assumptions prove $Sa\ne0$.
\end{lemma}
\begin{proof}
For the constant-zero update put $Z(n)=F_{v.0}(n)$. The two computation
schemas give $Z(0)=1$ and $Z(Sa)=0$. An equality $Sa=0$ would therefore
give $0=1$, contradicting \eqref{eq:ac-sep} by symmetry. No range theorem
for $Z$ and no constructor coverage are used.
\end{proof}

\subsection{An explicit finite-table compiler}

For a constant update define
\[
 D(n,a):=F_{v.a}(n).
\]
The computation schemas give
\[
 D(0,a)=Sa,\qquad D(Sn,a)=a.
\]
For a Boolean term $b$, meaning that $b=0\lor b=1$ has been proved,
define the conditional increment
\[
 \operatorname{Inc}(b,a):=D(\epsilon(b,0),a).
\]
Reflection and separation prove $\epsilon(0,0)=1$ and
$\epsilon(1,0)=0$; in the latter case use the two-value range axiom and
rule out the alternative value $1$. Thus
\begin{equation}
 \operatorname{Inc}(0,a)=a,\qquad
 \operatorname{Inc}(1,a)=Sa.
 \label{eq:ac-inc}
\end{equation}
These equations do not cancel a successor.

Fix a positive ordinary integer $k$, and suppose that
$\bar0,\ldots,\bar k$ have already been proved pairwise distinct.
The following notation is an external finite expansion:
\[
 T_0(v):=0,\qquad
 T_{j+1}(v):=
 \operatorname{Inc}\bigl(\epsilon(v,\bar j),\,\cdot\,\bigr)^{j+1}
                      \bigl(T_j(v)\bigr)
 \quad(0\le j<k),\qquad
 h_k(v):=T_k(v).
 \tag{Finite cycle term}\label{eq:ac-cycle-term}
\]
Here the exponent $j+1$ means exactly $j+1$ written compositions of the
displayed unary term template, not an object-level recursion or a new
iteration operation. Consequently \eqref{eq:ac-cycle-term} is a canonical
finite term in $0,S,\epsilon,F$ after the abbreviations are expanded.

For a fixed $i<k$, substituting $v=\bar i$ makes its $i$th equality test
equal to $1$ and all other displayed tests equal to $0$. Exactly $i+1$
conditional increments are then performed. For $v=\bar k$ all tests are
zero. Hence the source theory proves the finite table
\begin{equation}
 h_k(\bar i)=\overline{i+1}\ (0\le i<k),\qquad
 h_k(\bar k)=0.
 \label{eq:ac-table}
\end{equation}
It also proves, for arbitrary $v$,
\begin{equation}
 h_k(v)=\bar0\ \lor\cdots\lor\ h_k(v)=\bar k.
 \label{eq:ac-h-range}
\end{equation}
Indeed, make the finitely many case distinctions supplied by
\eqref{eq:ac-range}. If some displayed test is $1$, reflection identifies
$v$ with its corresponding numeral; distinctness makes all the other
tests zero, and the preceding computation applies. If all tests are zero,
the output is zero. This is a finite use of intuitionistic disjunction
elimination. It assumes no excluded middle for quantified formulas.

This construction also describes the exact finite-table capacity used
here: a numeral-valued table is a finite sum of disjoint Boolean
indicators, with each unit contribution implemented by
\eqref{eq:ac-inc}. It does not provide a selector between arbitrary
numerical values or an object-level evaluator of variable-length tables.

\begin{lemma}[The finite update has no $k$-periodic point]
\label{lem:ac-finite-update}
Under the displayed numeral-distinctness assumption, the source proves
\[
                       h_k^{\,k}(w)\ne w
                       \qquad\text{for every }w.
\]
Again, the superscript denotes a fixed finite composition.
\end{lemma}
\begin{proof}
If $h_k^{\,k}(w)=w$, then $k\ge1$ and the range theorem
\eqref{eq:ac-h-range} imply that $w$ is one of
$\bar0,\ldots,\bar k$. On these $k+1$ values, the fixed table
\eqref{eq:ac-table} is one cycle. Starting at any of them and applying
that table $k$ times ends at a different displayed numeral. Each case
therefore contradicts the assumed distinctness. The modular description
of this finite table is explanatory: all cases are finite computations
in the source theory.
\end{proof}

\begin{theorem}[Every fixed successor cycle is impossible]
\label{thm:ac-aperiodicity}
For each fixed ordinary integer $k\ge1$, assumptions
\eqref{eq:ac-F0}--\eqref{eq:ac-reflection} prove
\[
                           \forall a\;S^k a\ne a.
\]
No instance of numerical induction and no successor-injectivity axiom is
used in this family of proofs.
\end{theorem}
\begin{proof}
We construct these finite proofs by external induction on $k$.
Suppose all positive cycle lengths less than $k$ have been excluded.
The numerals $\bar0,\ldots,\bar k$ are pairwise distinct: equality of
$\bar0$ with a positive numeral contradicts Lemma~\ref{lem:ac-zero};
and equality $\bar i=\bar j$ with $1\le i<j\le k$ would give
\[
                   S^{j-i}(\bar i)=\bar i,
\]
contradicting an earlier theorem because $1\le j-i<k$. For $k=1$ the
second case is empty, so the construction starts from separation alone.

Construct the term $h_k$ of \eqref{eq:ac-cycle-term}. Suppose
$S^k a=a$, and put $w=F_{h_k}(a)$. Exactly $k$ applications of the
successor computation schema, followed by equality congruence, give
\[
 w=F_{h_k}(a)=F_{h_k}(S^k a)=h_k^{\,k}(F_{h_k}(a))=h_k^{\,k}(w).
\]
This contradicts Lemma~\ref{lem:ac-finite-update}. Discharge the
assumption and universally generalize $a$.

The induction in the first sentence is in the metatheory and produces
one finite source proof for each fixed $k$. It is not a numerical
induction instance of the theory being simplified.
\end{proof}

\begin{corollary}
Every fixed pair of distinct numerals is provably distinct, and, for
fixed $0\le i<j$, the source proves $S^i a\ne S^j a$ for arbitrary $a$.
\end{corollary}
\begin{proof}
The numeral assertion is part of the construction. For the second
assertion use Theorem~\ref{thm:ac-aperiodicity} with the element $S^i a$
and cycle length $j-i$.
\end{proof}

\subsection{Precisely what remains open}

The theorem excludes cycles along one successor chain. It does not infer
$a=b$ from $Sa=Sb$ for two arbitrary elements. A directed successor
structure can have two different predecessor chains merge and still have
no finite cycle. This observation identifies the gap; it is not a
countermodel to the full scalar-iteration theory.

The finite-table construction has a different fixed term $h_k$ for each
external $k$. It does not supply a term for cycles of arbitrary internal
length. The latter would require a uniform arithmetic construction with
proved internal invariants, and could plausibly support a cancellation
argument. It cannot be obtained merely by quantifying over the finite
proofs above: the source has neither a uniform total evaluator nor an
arbitrary-initializer iterator among its assumptions.

In particular, the existing scalar self-seeded completion uses injectivity
when its capped trajectory proves
$S(V(m,S^3m))=S^3m$ and concludes $V(m,S^3m)=S^2m$. The aperiodicity
theorem does not justify this cancellation. Deleting injectivity from that
completion on the basis of the present theorem would therefore be
unproved. The result is a new independently useful consequence and a
more precise research target, not a completed smaller axiomatization.

% Includable fragment. The construction was proposed by the root research
% agent and independently checked here. Requires the F/D/epsilon setting
% and theorem environments from the parent paper.
\section{Coverage and a right predecessor without successor injectivity}
\label{sec:right-predecessor}

There is a stronger consequence of the self-seeded equations than finite
aperiodicity. After reinstating contextual induction on literal equations,
the source can prove numerical coverage and construct a finite-case
right predecessor, still without successor injectivity. The word
``right'' is deliberate: the result is $S(P(Sn))=Sn$, not $P(Sn)=n$.

Use the computation, equality-test, and separation axioms stated in
Section~\ref{sec:finite-aperiodicity}, together with the contextual
literal-equation induction rule. Write $D(n,a)=F_{v.a}(n)$ as before.
All functions introduced below, except the explicitly graph-defined $P$,
are finite source terms.

\subsection{Quantifier-free induction is available without cancellation}

For Boolean values define
\[
 \neg_B b:=\epsilon(b,0),\qquad
 b\land_B c:=\epsilon\bigl(D(b,\epsilon(c,0)),0\bigr),\qquad
 b\lor_B c:=\neg_B(\neg_B b\land_B\neg_B c).
\]
Their truth tables use only $D(0,a)=Sa$, $D(1,a)=a$, equality reflection,
and $0\ne1$. In the conjunction table, the sole non-Boolean intermediate
value is $S1$; its distinction from $0$ follows from zero separation,
without requiring $S1\ne1$ or cancelling a successor.

External recursion on a fixed quantifier-free formula $A$ therefore
constructs a Boolean term $\chi_A$ with $\chi_A=1\leftrightarrow A$.
Implication is represented by $\neg_B\chi_A\lor_B\chi_B$; this is valid
intuitionistically because the finite equality characteristics make the
quantifier-free antecedent decidable. Literal-equation induction on
$\chi_A=1$ consequently yields contextual quantifier-free induction.
No successor-injectivity instance is used in this conversion.

\subsection{A clock that stops one step before its target}

For a passive parameter $a$, define
\[
 b_a(v):=\epsilon(v,0)\lor_B\epsilon(Sv,a),\qquad
 h_a(v):=D(b_a(v),v),\qquad
 V(j,a):=F_{h_a}(j),\qquad W(a):=V(a,a).
\]
The Boolean guard gives the exact two cases
\[
 h_a(v)=v\quad\text{if }v=0\text{ or }Sv=a,\qquad
 h_a(v)=Sv\quad\text{otherwise}.
\]
In particular $h_a(0)=0$, so the self-seeded base is
\[
 V(0,a)=1,\qquad V(Sj,a)=h_a(V(j,a)).
\]

\begin{lemma}[Capped-clock invariant]
\label{lem:rp-clock}
The source proves, for all $j,a$,
\[
                       V(j,a)=Sj\quad\lor\quad S(V(j,a))=a.
\]
\end{lemma}
\begin{proof}
Use quantifier-free induction on $j$, keeping $a$ fixed. The base is
$V(0,a)=1=S0$. For the step put $v=V(j,a)$. If the second alternative
$Sv=a$ holds, the update fixes $v$, preserving that alternative. Otherwise
the old first alternative gives $v=Sj$, so $v\ne0$. The guard is then
false, and the update gives $Sv=S(Sj)$, the first alternative at $Sj$.
Equivalently, one may split the Boolean guard first; a successful guard
in the first alternative cannot be caused by $v=0$ and therefore gives
$Sv=a$. Every case uses only the displayed induction hypothesis and
zero/successor separation.
\end{proof}

\begin{lemma}[Right predecessor on double successors]
\label{lem:rp-double}
For every $m$, the source proves
\[
                         S(W(S^2m))=S^2m.
\]
\end{lemma}
\begin{proof}
First use Lemma~\ref{lem:rp-clock} with its passive parameter still
independent, and only then specialize $a=S^2m$ and $j=m$. In its first
alternative $V(m,a)=Sm$; its second alternative already states
$S(V(m,a))=a$. Thus both alternatives prove the latter equation.
The guard therefore fixes $V(m,a)$. Two explicit successor computation
steps give
\[
 V(S^2m,a)=h_a(h_a(V(m,a)))=V(m,a).
\]
Since $a=S^2m$, the left side is $W(S^2m)$. Applying $S$ establishes
the conclusion. No successor is cancelled.
\end{proof}

\subsection{A finite graph and a legitimate induction instance}

Define $P$ by the finite quantifier-free graph
\[
\begin{split}
 \Phi_P(n,p):={}&(n=0\land p=0)\ \lor\ (n=1\land p=0)\\
 &{}\lor\ (n\ne0\land n\ne1\land p=W(n)).
\end{split}
\]
Equality decidability and $0\ne1$ prove this graph total and single
valued. Introduce $P$ in its conservative definitional extension.
This use of $P$ does not enlarge the primitive iteration schema to new
updates containing $P$.

For the induction needed next, finite-case elimination supplies more than
mere conservativity. Each term with finitely many $P$ occurrences has a
finite partition into source quantifier-free guards with source value
terms: refine cases for its arguments, and then use the three displayed
cases of $\Phi_P$. Repeating this finite expansion translates any
quantifier-free formula of the extension into an equivalent source
quantifier-free formula. Consequently contextual quantifier-free induction
is admissible in this extension. No general principle that definitional
extensions preserve restricted induction is being assumed.

\begin{theorem}[Coverage and a right predecessor are derived]
\label{thm:rp-coverage}
The source, through the finite graph above, proves
\[
 P0=P1=0,\qquad
 n=0\ \lor\ n=S(Pn),\qquad
 S(P(Sn))=Sn.
\]
It consequently proves $n=0\lor\exists p\;n=Sp$, without a coverage
axiom or successor-injectivity axiom.
\end{theorem}
\begin{proof}
The two boundary equations are immediate from the graph. Use the
legitimate quantifier-free induction just described on
$n=0\lor n=S(Pn)$.

At zero the first alternative holds. For the successor step, if the old
first alternative $n=0$ holds, then $Sn=1=S(P1)$. In the old second
alternative, $Sn=S^2(Pn)$, so Lemma~\ref{lem:rp-double} gives
$S(W(Sn))=Sn$. The input $Sn$ is nonzero. If $Sn=1$, its graph value
is zero and $S(P(Sn))=1=Sn$. If $Sn\ne1$, its graph value is $W(Sn)$,
and the same conclusion follows from the lemma. This proves the
successor instance of the induction formula in every case.

Apply the resulting coverage formula to $Sn$ and exclude its zero
alternative to obtain $S(P(Sn))=Sn$. Existential introduction gives the
ordinary coverage statement. All uses of $P$ can be removed by the
finite graph translation.
\end{proof}

\paragraph{The remaining cancellation problem.}
The theorem chooses one predecessor of each nonzero element. It does not
prove that every predecessor of that element is the chosen one. Equating
$Sa$ and $Sb$ permits applying $P$, but only gives $P(Sa)=P(Sb)$; the
missing statements are $P(Sa)=a$ and $P(Sb)=b$. Thus this construction
removes coverage from the no-injectivity precursor theory, but does not
yet remove the injectivity axiom from the existing complete iPRA
presentation.

\begin{lemma}[An exact form of the remaining closure problem]
\label{lem:rp-retraction}
Put $T(n):=P(Sn)$ and $C(n):\!\iff T(n)=n$ in the finite-case
definitional extension. Then
\[
 S(T(n))=Sn,\qquad T(T(n))=T(n),\qquad T(0)=0.
\]
Over the precursor theory, successor injectivity is equivalent to
\[
                         \forall n\;(C(n)\to C(Sn)).
\]
\end{lemma}
\begin{proof}
The first equation is Theorem~\ref{thm:rp-coverage}. Applying $P$ to it
gives the second; the third is $P1=0$. Thus $T$ is a retraction onto
its fixed points, with $0$ among them.

If these fixed points are successor closed, the already justified
quantifier-free induction on $C(n)$ gives $T(n)=n$ for every $n$.
From $Sa=Sb$, application of $P$ then gives $a=b$. Conversely,
successor injectivity applied to $S(T(n))=Sn$ gives $T(n)=n$ for
every $n$, hence the displayed closure condition.
\end{proof}

\paragraph{Why a direct comparison of the clocks is insufficient.}
One might try to establish this closure by comparing the updates with
targets $Sn$ and $S^2n$. For $v\ne0$, the expression $S(h_{Sn}(v))$
uses its stopping branch when $Sv=Sn$, whereas $h_{S^2n}(Sv)$ uses
its stopping branch when $S^2v=S^2n$. The putative intertwining
identity therefore needs
\[
                         S^2v=S^2n\ \Longrightarrow\ Sv=Sn
\]
at least for every state visited by the clocks. In the contrary case
the two sides are $S^2v$ and $Sv$, which are distinct by fixed-point
exclusion. The hypothesis $C(n)$ specifies the chosen predecessor of
$Sn$; no proved invariant turns it into uniqueness at $S^2n$.
This identifies a gap in that attempted proof, not an independence
result.

\section{Validation, reproducibility, and further research}
\label{sec:validation}

The composition suite checks 596 complete-process expectations and ten
targeted mutation observations. It covers direct computations, repaired
zero computations, nested binders, contextual induction, transport,
scope failures, false conclusions, scalar-major aliases and malformed
JSON, under ordinary Python and \texttt{-O}. The mutation controls exercise
major verification, equation-only induction, closed-goal scope, genuine
successor construction and lexical-floor transmission. Finite testing is
evidence about these implementations, not a proof of the universal
metatheorems.

One historical raw alias has a documented interface difference: the
545-byte guarded comparator rejects a null-major packet when its affine
dispatch attempts \texttt{None*3}; its translated empty-list major is a
valid accepted zero-instance packet. Canonical completeness does not
promise preservation of every redundant raw spelling.

The independent arithmetic diagnostics evaluate the finite-table and
capped-clock constructions on standard finite inputs. They are kept
separate from kernel certificates and never described as proving the
claims in nonstandard models. The re-examined progression-binder prototype
is outside the selected checker comparison: its theory is already present
in the source corpus, and its 576-byte implementation is larger. It is not
a newly proposed replacement for the measured calculus.

Run \texttt{python3 reproduce.py} from a fresh extraction to remeasure
the selected modules and frozen corpus, run the public checker suites,
and repeat the arithmetic diagnostics. Run \texttt{python3 build.py}
to rebuild the paper with PDFLaTeX. No network access, external proof
oracle, or shell escape is required. The standalone expanded \TeX\ source
can also be compiled on its own. Rebuilding can change PDF metadata, so
the delivered hash manifest should be checked before rebuilding.

\paragraph{Further research questions.}
The most concrete remaining mathematical question is whether the
self-seeded precursor can upgrade its selected right predecessor to a
left predecessor. A positive proof of $P(Sn)=n$, or any other internal
left inverse, would remove the injectivity assumption from the known
completion. A negative answer needs a model of the complete precursor,
not merely an acyclic noninjective successor graph.

The finite-cycle proof uses a different fixed term for each external
cycle length. A uniform internal finite-cycle construction could provide
stronger cancellation information, but cannot be obtained by silently
adding an evaluator for all descriptions. The existing all-starts theorem
already combines the progression binder's Step and Base interfaces. The
remaining question there is whether that fusion admits a smaller complete
implementation; re-proving the fusion alone would repeat the prior work.
Finally, a compact Lean or Rocq formalization of the sorted opener,
lexical scope invariant, arithmetic bootstrap and both interpretation
composites would provide much stronger assurance than the present
written proofs and finite tests.


\clearpage
\appendix
\section*{Attribution and reading the appendices}
\addcontentsline{toc}{section}{Attribution and reading the appendices}
The following mathematical and representation proofs are reproduced from
the supplied repository, with their source identities preserved in
\file{evidence/copied_sources.json}. They are inherited dependencies, not
new claims of this investigation. Words such as ``new'' within a reproduced
section refer to that section's original development. Historical packet
interfaces in the arity appendix are antecedents of the current one;
Section~\ref{sec:composition} specifies the selected 529/537-byte interface
and its exact changes. No historical source-size claim overrides the
current measured comparison.

The first five appendices supply the native rule inventory, ordinary
interpretations, scalar predecessor, actual-equality reflection, and
general primitive-recursion compiler. The last two supply the sorted
array/opening proof and the neutral-major invariant. The code is small
because it represents these rules compactly, not because these
mathematical dependencies have been erased.

% New mathematical result: complementary introduction choices.
% Requires amsmath/amssymb/amsthm and theorem, lemma, proposition environments.
% The inherited arithmetic-to-iPRA bridge must be supplied separately.

\section{Complementary introductions with one local assumption}
\label{sec:choice-introduction}

The old paired introduction retains both the zero instance and an
induction instance when its matrix is an equation. Both together are
unnecessary. A certificate may choose either one, provided the induction
choice is restricted to a literal equation. This changes neither the
formula language nor the meaning of a formula. The omitted combination
is recovered by a short nested introduction--elimination construction.

\subsection{Complete native syntax and rule inventory}

The scalar numerical language is
\[
 t::=x\mid0\mid St\mid I(n,a;u.h).
\]
Only $h$ binds the state variable $u$; $n$ and $a$ do not.
The update may contain fixed numerical parameters. The computation
axioms in every retained context are
\[
 I(0,a;u.h)=a,\qquad
 I(Sn,a;u.h)=h[I(n,a;u.h)/u].
 \tag{I0, IS}\label{eq:choice-computations}
\]
The argument also applies unchanged to the full fixed-description PR
language with its projection, composition, and recursion equations.
No numerical variable ranges over function descriptions, and no uniform
evaluator is postulated.

Native formulas and reusable context clauses are
\[
 A::=(s=t)\mid\mathsf D_{x,y}A,\qquad
 \langle A;y.B(y)\rangle,\quad y\notin\operatorname{FV}(A).
\]
Bound variables are renamed before capture-free substitution.
The two variables $x,y$ in a paired formula are distinct.
An ordinary assumption $F$ is the constant clause
$\langle H;v.F\rangle$, where $v$ is absent from $F$ and $H$ is the
fixed computation theorem $I(0,0;u.0)=0$. Thus ordinary hypothesis
use is a clause Use preceded by this computation proof.

There are computation axioms and exactly the following inference
patterns:
\[
\frac{\Gamma\vdash A}{\Gamma\vdash B(t)}
\quad\mathsf{Use}
\quad(\langle A;y.B(y)\rangle\in\Gamma),
\]
\[
\frac{\Gamma\vdash\mathsf D_{x,y}B(x,y)\qquad
      \Gamma,\langle B(a,0);y.B(a,y)\rangle\vdash F}
     {\Gamma\vdash F}
\quad\mathsf{DE},
\qquad
\frac{\Gamma\vdash s=t\qquad\Gamma\vdash F(s)}
     {\Gamma\vdash F(t)}
\quad\mathsf{Tr},
\]
\[
\frac{\Gamma,B(t,\eta_\varepsilon(z))\vdash B(t,Sz)}
     {\Gamma\vdash\mathsf D_{x,y}B(x,y)}
\quad\mathsf{DI}_{\varepsilon},
\qquad
\eta_0(z)=0,\quad\eta_1(z)=z.
\tag{Choice DI}\label{eq:choice-di}
\]
Here $\varepsilon$ is a certificate choice in $\{0,1\}$, and
$\varepsilon=1$ is permitted only when $B$ has a literal equation
root. An equation beneath a binder is not a literal equation.
Call $\varepsilon=0$ the \emph{case choice} and $\varepsilon=1$
the \emph{progress choice}. Both retain exactly one ordinary
formula assumption. Compound matrices have only the case choice.

In DI the eigenparameter $z$ is fresh for $\Gamma$, the witness
$t$, and every fixed parameter of $B$. In DE the eigenparameter
$a$ is fresh for $\Gamma$, the conclusion $F$, and every fixed
parameter of $B$. In transport, $F(v)$ is any fixed native
formula template. Displayed local assumptions and clauses are
discharged. There are no primitive reflexivity, coverage,
separation, ordinary logical, or additional induction rules.

We call this calculus $\mathcal C$. The original comparison
calculus $\mathcal N$ has identical syntax, computation,
Use, DE, and transport, but its introduction is
\[
\frac{\Gamma,B(t,0),[B(t,z)]_{\mathrm{at}}\vdash B(t,Sz)}
     {\Gamma\vdash\mathsf D_{x,y}B(x,y)}.
\tag{Old DI}\label{eq:choice-old-di}
\]
The bracket supplies $B(t,z)$ exactly at a literal equation.

\subsection{Exact equivalence and the unused-witness construction}

\begin{lemma}[Structural admissibility]
\label{lem:choice-structural}
Weakening, exchange, contraction, renaming, ordinary-assumption cut,
and capture-free numerical substitution are admissible in $\mathcal C$.
\end{lemma}
\begin{proof}
Induct on derivations. Rename local eigenparameters before inserting
proofs or terms. Numerical substitution preserves a formula's root
constructor, so an allowed progress choice remains allowed.
The selected choice does not change under substitution or cut.
For ordinary-assumption cut, replace each use of its constant clause
by the supplied proof. Computation schemas are closed under
capture-free numerical substitution. These facts verify every case.
\end{proof}

\begin{theorem}[Identity-language equivalence]
\label{thm:choice-equivalence}
For every retained clause context $\Gamma$ and native formula $F$,
\[
 \Gamma\vdash_{\mathcal C}F
 \quad\Longleftrightarrow\quad
 \Gamma\vdash_{\mathcal N}F.
\]
The transformations are effective on finite derivations and leave
all numerical and formula syntax unchanged.
\end{theorem}
\begin{proof}
Each new introduction is an instance of the old one after weakening
by the omitted local assumption. For the case choice at an equation,
add the induction instance. For the progress choice, add the zero
instance. Compound case introduction is already the old rule.
Thus every new derivation translates to an old one.

For the reverse direction, transform an old derivation recursively.
Only its atomic DI case needs a construction. Suppose its matrix is
the equation $E(x,y)$ and its transformed premise is
\[
 \pi:\quad\Gamma,E(t,0),E(t,z)\vdash E(t,Sz).
\]
The outer new inference will be case DI with witness $t$. It is
therefore enough to prove $E(t,Sz)$ under
$\Gamma,E(t,0)$.

Rename the eigenvariable in $\pi$ to a fresh $v$. Choose a new
variable $u$ absent from $E(t,y)$ and all retained data. Under
$\Gamma,E(t,0)$, progress DI with witness zero introduces
\[
             \mathsf D_{u,y}E(t,y)
\]
from exactly the renamed premise $\pi[v/z]$.
The bound coordinate $u$ is deliberately unused.

Apply DE to this paired formula. Because $u$ is unused, the opened
clause is exactly
\[
               \langle E(t,0);y.E(t,y)\rangle;
\]
it contains no new existential witness. The retained ordinary
assumption $E(t,0)$ supplies its antecedent. Use at $Sz$ gives
$E(t,Sz)$, and DE closes. Finally the outer case DI introduces
the original $\mathsf D_{x,y}E(x,y)$ with witness $t$.
The original premise proof $\pi$ has been inserted once.

An old compound introduction already is the new case choice.
All other rule cases are unchanged. This completes the induction
on the original derivation.
\end{proof}

The unused inner coordinate is essential. Introducing
$\mathsf D_{x,y}E(x,y)$ internally and then eliminating it would
give an arbitrary fresh witness, not necessarily the original
$t$. The construction instead keeps $t$ as a fixed parameter
of $E(t,y)$ and binds only the unused coordinate $u$.

\begin{proposition}[Linear bound on inference nodes]
\label{prop:choice-node-bound}
Count each computation axiom, Use, DI, DE, and transport as one
inference node, expanding ordinary-assumption use as a computation
of $H$ followed by Use. If an old derivation has $N$ nodes and
$r$ atomic DI nodes, the displayed transformation has exactly
\[
                   N+5r\leq 6N
\]
nodes.
\end{proposition}
\begin{proof}
Each atomic old DI node is replaced by six nodes: outer case DI,
inner progress DI, DE, Use of the opened clause, Use of the
retained base assumption, and the computation of $H$ needed for
that last ordinary-assumption use. Its old premise subtree
occurs once. All other inference nodes are copied once.
\end{proof}

This counts inference-tree nodes. It is not a bound on serialized
JSON bytes, formula sizes, checker runtime, or the number of
consumed packet cells when hypothesis closure is implicit.
In the formula-context Q/D presentation, the corresponding
replacement has outer case DI, inner progress DI, DE, QE,
and two Hyp nodes, giving the same six-node count.

\subsection{Induction and numerical cases separate cleanly}

\begin{lemma}[Contextual equation induction]
\label{lem:choice-ei}
The progress choice, DE, and Use derive
\[
\frac{\Gamma\vdash E(0)\qquad
      \Gamma,E(z)\vdash E(Sz)}
     {\Gamma\vdash E(t)}
\tag{EI}\label{eq:choice-ei}
\]
for every equation $E$, with the usual freshness condition.
\end{lemma}
\begin{proof}
Choose an unused witness coordinate $u$. Progress DI introduces
$\mathsf D_{u,y}E(y)$ from the given step. DE opens the clause
$\langle E(0);y.E(y)\rangle$. Supply its antecedent using the given
base proof and use the clause at $t$. Close DE.
\end{proof}

\begin{lemma}[Cases, equality, and binder congruence]
\label{lem:choice-preinduction}
The case choice derives, for arbitrary native formulas,
\[
\frac{\Gamma\vdash F(0)\qquad\Gamma\vdash F(Sz)}
     {\Gamma\vdash F(t)}
\tag{NC}\label{eq:choice-nc}
\]
with the usual freshness condition. Equality reflexivity,
symmetry, transitivity, numerical congruence, and uniform
matrix interderivability under $\mathsf D$ are derivable
without using the progress choice.
\end{lemma}
\begin{proof}
For cases, use case DI with an unused witness coordinate and
matrix $F(y)$, ignoring its retained base in the successor
premise. DE opens $\langle F(0);y.F(y)\rangle$; the given zero
proof and Use at $t$ finish.

For reflexivity put $r=I(0,t;u.0)$. Use the computation equation
$r=t$ as both transport premises with motive $v=t$ to obtain
$t=t$. The usual explicit equality motives give symmetry,
transitivity, and numerical congruence.

For binder congruence, open the old paired formula at a fresh
witness $a$ and introduce the new one by case DI at $a$.
Convert its retained base to the old base, use the opened
old clause at the successor, and convert back. Ordinary cut
closes the construction, and the reverse conversion gives
the converse. No induction is used.
\end{proof}

The primitive progress choice is still restricted to equations.
Allowing it at arbitrary compound matrices would, by the
argument for \eqref{eq:choice-ei}, supply induction for
arbitrary formulas and change the intended theory.

\subsection{The entire ordinary logical interface}
\label{subsec:choice-logic}

Here is a predecessor-free reconstruction of the ordinary logical
operations. It makes explicit that the case choice supplies the
logical interface independently of progress induction.

Let $C(v,a,b):=I(v,a;u.b)$ with $u$ unused. Computation gives
$C(0,a,b)=a$ and $C(Sv,a,b)=b$. Define the gate on fixed finite
formula syntax by
\[
 G_v(s=t):=(C(v,0,s)=C(v,0,t)),\qquad
 G_v(\mathsf D_{x,y}B):=\mathsf D_{x,y}G_vB,
\]
renaming bound variables away from the control. Equality,
binder congruence, and external structural induction give
\[
 \vdash G_0A,\qquad G_{Sv}A\dashv\vdash A,\qquad
 A\dashv\vdash\mathsf D_{u,y}G_yA
 \quad(u,y\notin\operatorname{FV}(A)).
\tag{Gate and padding}\label{eq:choice-gate}
\]
For padding, case DI proves the forward direction. DE, the
zero-gate theorem, and Use at $1$ prove the reverse direction.

Define a selector between two arbitrary formulas by
\[
 M_v(s=t,u=w):=(C(v,s,u)=C(v,t,w))
\]
and, with aligned fresh paired variables, by
\[
\begin{aligned}
 M_v(\mathsf D_{x,y}A,\mathsf D_{x,y}B)
   &:=\mathsf D_{x,y}M_v(A,B),\\
 M_v(E,\mathsf D_{x,y}B)
   &:=\mathsf D_{x,y}M_v(G_yE,B),\\
 M_v(\mathsf D_{x,y}A,E)
   &:=\mathsf D_{x,y}M_v(A,G_yE),
\end{aligned}
\]
where $E$ is an equation. Every call decreases maximum input
binder depth. Padding and congruence prove
$M_0(A,B)\dashv\vdash A$ and
$M_{Sv}(A,B)\dashv\vdash B$.

Put $\delta(v)=C(v,1,0)$ and $p(z)=I(z,0;u.\delta(u))$.
Only the fixed computations $p(1)=1$ and $p(2)=0$ are needed.
The full logical macro inventory is
\[
\begin{aligned}
 A\Rightarrow_*B&:=\mathsf D_{u,z}M_z(A,B),\\
 A\wedge_*B&:=\mathsf D_{u,z}G_zM_{p(z)}(A,B),\\
 \exists_*x\,A(x)&:=\mathsf D_{x,y}G_yA(x),\\
 A\vee_*B&:=\exists_*x\,M_x(A,B),\\
 \forall_*x\,A(x)&:=
        A(0)\wedge_*\mathsf D_{u,z}G_zA(z),\\
 \bot_*&:=(0=1).
\end{aligned}
\tag{Logical macros}\label{eq:choice-macros}
\]
Newly bound variables are fresh; $u$ is unused.

Implication introduction obtains its assumption from the zero
selector and supplies the consequent at successors, then uses
case DI. Elimination opens its clause and uses it at $1$.
Existential introduction uses its witness and the successor
gate; elimination opens at a fresh witness, supplies the
zero-gate theorem, and uses the clause at $1$.

Disjunction introduction uses witnesses $0$ and $1$.
For elimination, discharge the two branch proofs as
$A\Rightarrow_*F$ and $B\Rightarrow_*F$. NC applied to
$M_z(A,B)\Rightarrow_*F$ gives this implication at every
control: the two endpoint cases follow from the branch
implications. Existential elimination provides $M_a(A,B)$
and the implication gives $F$.

Given $A$ and $B$, selector endpoints and NC give $M_v(A,B)$
at an arbitrary control. Substitute $v=p(Sz)$ and use the
successor gate for conjunction introduction. For projections,
open its zero-gated clause and specialize at $1$ and $2$;
the two fixed computations give $B$ and $A$, respectively.

These rules prove existential numerical coverage:
\[
     (x=0)\vee_*\exists_*v\,(x=Sv).
\tag{Coverage}\label{eq:choice-coverage}
\]
Use NC. Zero takes the left injection of reflexivity, and a
successor takes the right injection with the evident witness.
Coverage gives equality cases even when the retained context
contains $x$; its elimination adds $x=0$ or $x=Sv$ rather
than silently changing that retained context.

Universal introduction follows from a fresh-parameter proof
by substituting zero and successors and forming the two
conjuncts. For elimination, take its $A(0)$ conjunct and
open the paired conjunct. Its base is a zero-gate theorem.
Coverage of the target chooses the zero instance from
$A(0)$ or a successor instance from the opened clause
and gate removal.

Finally, under $0=1$, computation and transport give
$s=C(0,s,t)=C(1,s,t)=t$ for arbitrary numerical terms.
Every equation follows, and external induction on formula
syntax with case DI gives every native formula. This
supplies falsity elimination. No excluded middle for
quantified formulas has entered the construction.

\subsection{Whole-domain realizations and both formula composites}

For a standard first-order reading, let $T_{\rm eq}$ have
intuitionistic equality logic, the same numerical signature
and computation equations, numerical coverage, separation
$(0=1)\to\bot$, and contextual equation induction.
Its interpretation of the native syntax is
\[
 (s=t)^\circ=(s=t),\qquad
 (\mathsf D_{x,y}B)^\circ
   =\exists x\bigl(B^\circ(x,0)\to\forall y B^\circ(x,y)\bigr),
\]
\[
 \langle A;y.B\rangle^\circ
       =A^\circ\to\forall y B^\circ.
\]
The domain is the entire numerical domain and interpreted
equality is actual equality.

Case DI is sound in this realization by numerical coverage.
Progress DI is sound by contextual equation induction under
the conditional base. Use, DE, transport, and computation
have their displayed ordinary meanings. Conversely, the
native macros derive every ordinary logical rule,
\eqref{eq:choice-ei} derives equation induction, and
\eqref{eq:choice-coverage} derives coverage. Separation
translates to the identity implication
$(0=1)\Rightarrow_*(0=1)$.

For the essential paired formula composite,
\[
 \mathsf D_{x,y}B
 \dashv\vdash
 \exists_*x\bigl(B(x,0)\Rightarrow_*\forall_*y\,B(x,y)\bigr).
\tag{Paired composite}\label{eq:choice-paired-composite}
\]
Forwards, DE opens its clause at a fresh $a$; under $B(a,0)$
that clause supplies $B(a,y)$ at a fresh arbitrary $y$.
Introduce the universal, implication, and existential
macros and close DE. Backwards, eliminate the existential
at a fresh $a$. Introduce the paired formula with witness
$a$ by case DI: the retained base, supplied implication,
and universal elimination at $Sz$ give its successor premise.

The ordinary interpretations of the six logical macros
have their ordinary round trips. For conjunction, its true
base and the instances $1,2$ recover both operands;
the converse uses selector cases. The universal has all
successor instances in its paired conjunct and zero in
its first conjunct; coverage gives exactly the ordinary
universal. Implication, disjunction, and existential
quantification use their selector or gate endpoints.
Falsity uses the explicit ordinary separation bridge.
External induction on finite formulas, together with
\eqref{eq:choice-paired-composite}, gives both formula
composites. Clauses correspond in both directions to
$A\Rightarrow_*\forall_*yB$ by the derived logical rules.

The numerical composites are literal identities with
internally provable graph $x=y$. No standard-numeral
restriction, retract, quotient, or equality reflection
assumption is being inserted at this stage. The separately
supplied arithmetic bridge identifies $T_{\rm eq}$ with
full iPRA; the present change preserves that bridge
unchanged because Theorem~\ref{thm:choice-equivalence}
preserves every native sequent.

\subsection{What the theorem does and does not claim}

The new inference interface retains one ordinary assumption
at every paired introduction. It preserves the entire
consequence relation, not just equations or closed
theorems. Consequently the established full-iPRA
interpretations and equiconsistency transfer unchanged.

This result is different from retaining only progress
introduction and strengthening closure elimination.
Here the original elimination rule is retained, and
the case and progress introductions coexist. It is also
different from allowing arbitrary compound progress,
which would supply stronger induction.

A complete-checker comparison must still count the
certificate choice, its equation-root restriction,
substitution, context handling, computation, and the
public input/output interface. The mathematical theorem
alone is not a byte-count or minimality theorem.
The inference-node bound is distinct from source-size
measurements and from serialized certificate size.
The proof is a written universal proof; finite regression
tests do not constitute a Lean/Rocq formalization of it.

\section{Arithmetic strength and both ordinary interpretation composites}
\label{sec:ordinary-bridge}

We spell out the arithmetic bridge to avoid identifying standard numerical
agreement with bi-interpretability. First elaborate every fixed iterator
into an ordinary function symbol. At stage zero the symbols are $0,S$.
Given a fixed earlier-stage term $h(u,\vec p)$, introduce
$i_h(n,a,\vec p)$ with the two equations
\[
i_h(0,a,\vec p)=a,\qquad
i_h(Sn,a,\vec p)=h(i_h(n,a,\vec p),\vec p).
\]
Only finite descriptions occur; the domain contains no code for a
universal evaluator. Formal parameters are distinct when a description
is declared, while actual arguments can coincide.

Let $T_I$ be ordinary intuitionistic first-order logic with actual
equality and genuine falsity, these staged equations, contextual
induction for literal equations, and
\[
\forall x\,(x=0\vee\exists y\,x=Sy),
\qquad (0=S0)\to\bot.
\]
The second sentence is the falsity bridge. It must be stated: positive
numerical equations by themselves admit a singleton structure. Native
explosion from $0=S0$ and a translation of ordinary falsity to that
equation do not authorize silently dropping the bridge from $T_I$.

The native-to-ordinary translation is
\[
(s=t)^\circ=(s=t),\qquad
(\mathsf D_{x,y}B)^\circ=
\exists x\bigl(B^\circ(x,0)\to\forall y B^\circ(x,y)\bigr).
\]
A clause $\langle A;y.B\rangle$ becomes $A^\circ\to\forall yB^\circ$.
For the auxiliary closure used by the compact checkers,
$(\mathsf Q_yB)^\circ=B^\circ(0)\to\forall yB^\circ$.
The interpretation domain is all elements, defined by $x=x$;
equality is actual equality.

Case DI is sound because numerical coverage combines its retained base
with its successor proof. Progress DI is sound by literal-equation
induction under a temporarily assumed base. Use, DE and transport are
ordinary logical inferences under the translations. Conversely the
derived logical operations, equation induction, cases and explosion in
Section~\ref{sec:choice-introduction} recover the rules and axioms of
$T_I$. An auxiliary closure is interderivable with
$\mathsf D_{v,y}B(y)$ for unused $v$: in one direction use case DI with
the closure and its base; in the other eliminate the paired formula
and return its supplied closure. Thus formulas and context clauses have
the stated meanings in both presentations.

\subsection{The arithmetic dependency order}

The scalar update has no recursion-counter argument. We do not conceal
one in an encoding. The inherited internal proofs, reproduced in the
appendices, establish the required arithmetic in this order:
\begin{enumerate}
\item Using equation induction and numerical cases, construct a fixed
scalar predecessor term and prove $P0=0$ and $P(Sn)=n$.
Existential coverage then gives $x=0\vee x=S(Px)$: in the successor
case $x=Sy$, apply $P$ to obtain $Px=y$ and substitute.
\item Construct truncated subtraction and a numerical characteristic
of actual equality. Prove its reflection on the entire original domain.
\item Compile each fixed quantifier-free formula into its characteristic
equation and recover every contextual quantifier-free induction instance.
\item With that induction available, construct pairing and projections
by a bounded square-root search and compile every fixed general
primitive-recursive description into a scalar iterator term.
\end{enumerate}
No pairing or faithful equality detector is assumed to obtain
predecessor. No appeal to standard-model correctness replaces an
internal equality. Every induction fixes its passive parameters;
diagonal substitution occurs only into a previously proved identity.

For a target recursion
\[
f(0,\vec p)=g(\vec p),\qquad
f(Sn,\vec p)=h(n,f(n,\vec p),\vec p),
\]
let $\pi,L,H$ be the proved pairing and projections. The compiler uses
\[
V(n,\vec p)=I\bigl(n,\pi(0,g^*(\vec p));
u.\pi(SL(u),h^*(L(u),H(u),\vec p))\bigr),
\qquad f^*(n,\vec p)=H(V(n,\vec p)).
\]
Equation induction proves $L(V(n,\vec p))=n$; the two projection
identities then prove the target defining equations. The construction
is external recursion on each finite PR description, not an internal
quantification over arbitrary functions.

\begin{theorem}[Whole-domain bi-interpretability and equiconsistency]
The ordinary realization $T_I$ and full first-order iPRA are
bi-interpretable. In both theories $x=y$ defines the isomorphism from
the original domain to its double interpretation. The native Choice-DI
calculus is therefore a faithful proof presentation of a theory
bi-interpretable and equiconsistent with iPRA.
\end{theorem}
\begin{proof}
Interpret an iterator in iPRA by primitive recursion that ignores its
counter argument in the update. iPRA proves its equations, literal
equation induction, coverage and separation. In the reverse direction
use the scalar compiler. The predecessor and equality-reflection proofs
supply the successor axioms and every quantifier-free induction instance.
All ordinary logical rules are preserved.

For two iterators with the same seed and provably equal updates,
equation induction on the equality of their values proves uniqueness at
every counter, with all parameters fixed. The corresponding argument
applies to ordinary primitive recursion. Induct externally on a finite
description: its double translation satisfies its original defining
equations, so internal uniqueness identifies the two functions. This
also proves coherence when substitution changes the staged syntactic
description of an update.

The domain and equality are unchanged; hence the identity graph is
internally total, bijective and preserves every interpreted operation.
The displayed ordinary meanings of the logical macros and the paired
binder give both formula composites by external induction on finite
formula syntax. Clauses are similarly equivalent to their conditional
universal meanings. The explicit falsity bridge preserves contradiction
in both directions. Effective translations of finite proofs give
equiconsistency, without a claim that either theory proves its own
consistency.
\end{proof}

% Adapted with attribution from the pinned unparenthesized one-hole archive.
\section{Internal scalar predecessor bootstrap}
\label{sec:hybrid-bootstrap}

This scalar bootstrap is inherited from the supplied scalar development~\cite{erasure} and is reproduced to keep the scalar checker self-contained. It is not a new contribution of this paper. The proof uses a factor map fixing odd states and a saturated-iteration construction of predecessor.

The available arithmetic terms are $0$, $S$, and
\[
 I(n,a;u.h),\qquad
 I(0,a;u.h)=a,\qquad
 I(Sn,a;u.h)=h[I(n,a;u.h)/u].
\]
The update binds only $u$.  All other variables in it are fixed numerical
parameters.  We use equality reasoning, the already reconstructed
zero/successor cases, and induction on a single literal equation.
Every induction below keeps its other parameters fixed.  Predecessor,
successor injectivity, pairing, and induction on arbitrary
quantifier-free formulas are not assumptions of the argument.

\subsection{Elementary terms and two finite cycles}

Write $1=S0$ and use larger numerals as successor abbreviations.  Define
\[
 C(n,a,b)=I(n,a;u.b),\qquad
 Z(n)=C(n,1,0),\qquad
 a+b=I(b,a;u.Su).
\]
The update variable of $C$ is absent from its body.  Consequently
\[
 C(0,a,b)=a,\qquad C(Sn,a,b)=b.
\]
Numerical cases show that $Z(n)$ is $0$ or $1$.
The usual commutative-monoid equations for addition follow by literal
equation induction.  For example, $0+b=b$ and $Sa+b=S(a+b)$ are
proved by induction on $b$; these yield commutativity by a further
equation induction, and associativity follows by induction on the last
summand.  We use these equations to reorder finite sums.  The notation
$2a$ means $a+a$ and $3a$ means $a+a+a$; no multiplication operation
is needed for this bootstrap.

Put
\[
 N(n)=I(n,0;u.Z(u)),\qquad
 j(u)=Z(u)+2N(u),\qquad
 M(n)=I(n,0;u.j(u)).
\]
The range of $N$ is contained in $\{0,1\}$: its zero value is $0$,
and its successor value is a value of $Z$.  This uses numerical cases
on its counter, not induction on a disjunction.  Hence
\begin{equation}
 N(n)+Z(N(n))=1,\qquad
 Z(Z(N(n)))=N(n),\qquad N(n+2)=N(n).
 \label{eq:hybrid-parity}
\end{equation}
Here the last equation follows from two successor computations and
the preceding equation.  Literal-equation inductions give
\begin{equation}
 N(2n)=0,\qquad N(2n+1)=1,\qquad N(3n)=N(n).
 \label{eq:hybrid-parity-scaling}
\end{equation}
For the first equation the successor step adds two to the argument of
$N$; the second follows by one computation.  In the third, increasing
$n$ adds three to the left argument and therefore applies $Z$ once to
both sides of the induction equation.

Finite computations now give
\[
 j(0)=1,\qquad j(1)=2,\qquad j(2)=0.
\]
The proof below needs just these three values and
$M(Sn)=j(M(n))$; it does not assume a division algorithm or a theorem
characterizing $M$ as remainder modulo three.

\subsection{A ramp with a global scaling identity}

Define
\[
 G(n)=I(n,0;u.u+M(u)+1).
\]
The central invariant is
\begin{equation}
 M(G(n))=N(n).
 \label{eq:hybrid-ramp-invariant}
\end{equation}
Its zero case is immediate.  For its successor step, retain the
single induction equation and split the Boolean value $N(n)$.
If it is $0$, the next $G$ state is $G(n)+1$; its $M$ value is
$j(0)=1$, equal to $N(Sn)$.  If it is $1$, the next $G$ state
is $G(n)+2$; two computations of $M$ change its value from $1$
to $2$ and then to $0$, again equal to $N(Sn)$.  This proves
the invariant by literal-equation induction.

It follows that
\[
 G(Sn)=G(n)+N(n)+1,\qquad G(n+2)=G(n)+3.
\]
Equation induction therefore yields
\begin{equation}
 G(2n)=3n.
 \label{eq:hybrid-ramp-double}
\end{equation}
A separate literal-equation induction proves
\begin{equation}
 G(3n)=3G(n)+N(n).
 \label{eq:hybrid-ramp-triple}
\end{equation}
For completeness, its successor calculation is
\begin{align*}
 G(3Sn)
 &=G(3n)+N(3n)+4
   =3G(n)+2N(n)+4,\\
 3G(Sn)+N(Sn)
 &=3G(n)+3N(n)+3+Z(N(n)).
\end{align*}
The expressions agree in both Boolean cases for $N(n)$.
The base is $0=0$.  In particular,
\eqref{eq:hybrid-ramp-triple} was proved internally; it was not
inferred from an intended floor-function interpretation of $G$.

Set
\[
 H(u)=C(N(u),G(u),u).
\]
The crucial choice is to retain $u$ when its parity is $1$.
The preceding equations give
\begin{equation}
 H(2u)=3u,\qquad H(1)=1,\qquad H(3u)=3H(u).
 \label{eq:hybrid-factor-map}
\end{equation}
To check the last identity, use $N(3u)=N(u)$ and split the Boolean
$N(u)$.  At $0$, both sides are $3G(u)$ by
\eqref{eq:hybrid-ramp-triple}; at $1$, both are $3u$.
This argument does not require that an arbitrary number has already
been represented as $2v$ or $2v+1$.

\subsection{A comparator from scalar iteration}

For every fixed update $h$, equation induction on $n$, keeping $m,a$
and all parameters fixed, proves the semigroup identity
\[
 I(m+n,a;h)=I(n,I(m,a;h);h).
\]
The base is the zero computation and the step applies $h$ to the
induction equation. This identity will also be used in the later
pairing construction.


For a fixed update $h$, equation induction on $y$, with its initializer
$a$ and all parameters of $h$ fixed, proves the rotation identity
\begin{equation}
 I(Sy,a;h)=I(y,h(a);h).
 \label{eq:hybrid-rotation}
\end{equation}
Both sides have base $h(a)$, and their successor equations apply the
same $h$ to the previous equality.  Here and below $I(y,a;h)$ is
an abbreviation for one fixed finite scalar term, not quantification
over functions.

By \eqref{eq:hybrid-factor-map}, another equation induction gives
\begin{equation}
 I(y,3a;H)=3I(y,a;H).
 \label{eq:hybrid-iterator-scaling}
\end{equation}
Its base is $3a=3a$.  In the step the fixed-parameter induction
equation rewrites the new left side as
\[
 H(3I(y,a;H))=3H(I(y,a;H)).
\]
The initializer parameter stays $a$ throughout this induction.
Equation induction also gives $I(y,1;H)=1$, since $H(1)=1$.

Define
\[
 E(x)=I(x,1;u.2u),\qquad
 L(x,y)=N\bigl(I(y,E(x);H)\bigr).
\]
Then the following entire recursion specification is provable:
\begin{equation}
 L(0,y)=1,\qquad L(Sx,0)=0,\qquad L(Sx,Sy)=L(x,y).
 \label{eq:hybrid-comparator}
\end{equation}
The first equation uses $E(0)=1$ and the fixed-point identity for
$H$.  The second uses $E(Sx)=2E(x)$ and $N(2a)=0$.
For the third,
\begin{align*}
 L(Sx,Sy)
 &=N\bigl(I(Sy,2E(x);H)\bigr)\\
 &=N\bigl(I(y,H(2E(x));H)\bigr)\\
 &=N\bigl(I(y,3E(x);H)\bigr)\\
 &=N\bigl(3I(y,E(x);H)\bigr)\\
 &=N\bigl(I(y,E(x);H)\bigr)=L(x,y).
\end{align*}
The successive steps use rotation, $H(2a)=3a$, iterator scaling,
and $N(3a)=N(a)$.  Thus $L$ has the comparison equations at all
arguments internally.  Its standard-model meaning as the indicator
of $x\le y$ is motivation, not a premise of the following argument.

\subsection{Saturated iteration gives predecessor}

We now need only the three equations
\eqref{eq:hybrid-comparator}.  Define
\[
 B_m(u)=C(L(Su,m),u,Su),\qquad
 W(t,m)=I(t,0;u.B_m(u)).
\]
The subscript in $B_m$ records a fixed parameter; $B_m$ is an open
term abbreviation, not a function variable.

Numerical cases on $e$ prove the selector equation
\[
 C(e,Sa,Sb)=S C(e,a,b).
\]
Equivalently one may use equation induction on $e$, whose induction
hypothesis is unused.  Therefore
\begin{align}
 B_{Sm}(Su)
 &=C(L(SSu,Sm),Su,SSu)\notag\\
 &=C(L(Su,m),Su,SSu)=S B_m(u),
 \label{eq:hybrid-conjugacy}\\
 B_{Sm}(0)&=C(L(S0,Sm),0,S0)=C(L(0,m),0,S0)=S0.
 \label{eq:hybrid-conjugacy-base}
\end{align}

Literal-equation induction on $t$, with $m$ fixed, proves
\begin{equation}
 W(St,Sm)=S W(t,m).
 \label{eq:hybrid-saturated-shift}
\end{equation}
The base is \eqref{eq:hybrid-conjugacy-base}.  The successor step is
\begin{align*}
 W(SSt,Sm)
 &=B_{Sm}(W(St,Sm))\\
 &=B_{Sm}(S W(t,m))\\
 &=S B_m(W(t,m))=S W(St,m).
\end{align*}
Only the second line uses the current induction hypothesis, at the
same fixed $m$.

Two subsequent equation inductions give
\begin{equation}
 L(Sn,n)=0,\qquad W(n,n)=n.
 \label{eq:hybrid-diagonals}
\end{equation}
The first uses the comparator recurrence.  The second uses
\eqref{eq:hybrid-saturated-shift} in its successor step.  The
substitution $t=m=n$ is made into that already proved universally
parameterized equation; no parameter of an active induction hypothesis
is changed.  Consequently
\begin{equation}
 W(Sn,n)=B_n(n)=C(L(Sn,n),n,Sn)=n.
 \label{eq:hybrid-extra-step}
\end{equation}

Finally define the finite scalar term
\begin{equation}
 P(m)=I\bigl(m,0;u.C(L(SSu,m),u,Su)\bigr).
 \label{eq:hybrid-predecessor-term}
\end{equation}
Its zero equation is immediate.  To justify its successor equation
without assuming a separate extensionality rule for iterator bodies,
put
\[
 V(t,m)=I\bigl(t,0;u.C(L(SSu,Sm),u,Su)\bigr).
\]
The comparator recurrence identifies this update pointwise with
$B_m(u)$.  Equation induction on $t$, with $m$ fixed, gives
$V(t,m)=W(t,m)$: their base values are $0$, and the successor step
uses the pointwise update equation and then the induction equation.
Now \eqref{eq:hybrid-extra-step} yields
\[
 \boxed{\qquad P(0)=0,\qquad P(Sm)=V(Sm,m)=W(Sm,m)=m.\qquad}
\]

\subsection{Dependency boundary}

All terms in this section are built from $0,S,I$ by finite substitution.
All inductions have a single equation as their invariant.  Numerical
cases and equality transport are the already reconstructed logical
resources.  Successor injectivity may now be obtained by applying $P$
to an assumed equation $Sx=Sy$.  Equality tests, reflection of
quantifier-free formulas, pairing, and the recovery of general
primitive recursion can then be developed in the established order.
They were not used in order to manufacture predecessor.

This argument is a written internal derivation in the stated theory.
The finite implementation checks do not constitute a proof-assistant
formalization of this entire metatheorem.

% Standalone section.  Requires only amsmath and amssymb in the parent document.
% No theorem environments, macros, or bibliography keys are assumed.
% All arithmetic claims in this section are proved here.

\section{From literal-equation induction to full intuitionistic PRA}
\label{sec:eic-completion}\label{app:arithmetic}

This section supplies an unconditional completion theorem. It is adapted
with attribution from the current forward-left draft~\cite{forward}. In the
scalar application used here, every function occurring in the proof is
already a scalar term: predecessor is supplied by
Appendix~\ref{sec:hybrid-bootstrap}, the selector is
$C(n,a,b)=I(n,a;u.b)$, addition is $x+y=I(y,x;u.Su)$, and
truncated subtraction is $d(x,y)=I(y,x;u.Pu)$. All other functions below
are finite compositions of these. Thus this proof does not presuppose
that arbitrary primitive recursion has already been compiled. The QFI
conclusion first applies to the scalar language; the general PR compiler
comes afterward.  Literal-equation
induction, together with the primitive-recursive defining equations,
predecessor coverage, and separation of $0$ and $1$, proves induction for every
quantifier-free formula.  In particular, the proof does not assume a faithful
numerical equality test, stability or decidability of equality, successor
uniqueness, or any commutation law for an iterator.  We construct the equality
test and prove its reflection property for arbitrary elements of the
numerical domain.

\subsection{The exact hypotheses}

Work in ordinary intuitionistic first-order logic with equality and genuine
falsity $\bot$.  The language has one numerical sort and a symbol for each
fixed, finite primitive-recursive description.  Distinct descriptions are
not identified merely because they compute the same standard function.
Let $\mathsf{PRE}$ contain the defining equations for zero, successor,
projections, composition, and primitive recursion.  In particular, choose
descriptions $P$ and $C$ with equations
\begin{align}
 P(0)&=0, & P(Sx)&=x, \label{eq:eic-predecessor}\\
 C(0,a,b)&=a, & C(Sx,a,b)&=b. \label{eq:eic-selector}
\end{align}
Write $1=S0$.  In addition to $\mathsf{PRE}$, assume
\begin{align}
 \mathsf{Cov}:&\qquad
       \forall x\,\bigl(x=0\lor x=S(Px)\bigr),
       \label{eq:eic-coverage}\\
 \mathsf{Sep}:&\qquad 0=1\longrightarrow\bot,
       \label{eq:eic-separation}
\end{align}
and induction for each \emph{literal equation} $E(n,\vec p)$:
\begin{equation}
 \forall\vec p\,
 \left[
   \left(E(0,\vec p)\land
       \forall n\bigl(E(n,\vec p)\longrightarrow E(Sn,\vec p)\bigr)
   \right)
   \longrightarrow\forall n\,E(n,\vec p)
 \right].
 \label{eq:eic-induction}
\end{equation}
Here the schematic expression $E$ is a single equality of numerical
terms; the displayed universal quantifiers range only over numbers.
Every description and every instance of the schema is finite external
syntax.

The axiom formulation \eqref{eq:eic-induction} and the usual contextual
induction rule are interchangeable.  For example, if
\[
 \Gamma\vdash E(0,\vec p),
 \qquad
 \Gamma,E(z,\vec p)\vdash E(Sz,\vec p),
\]
where $z$ is absent from $\Gamma$ and the passive parameters, discharge the
second premise's equation and universally introduce $z$.  The axiom then
gives $\Gamma\vdash\forall n E(n,\vec p)$, which can be instantiated at any
target term.  Conversely, assume the antecedent of the axiom and apply the
contextual rule with fresh variables.  This uses the ordinary intuitionistic
logical rules and permits arbitrary, including quantified, formulas in
$\Gamma$.

\subsection{Constructor separation and two elementary tests}

Define the sign description and the countdown description by
\begin{equation}
 \operatorname{sg}(u)=C(u,0,1),
 \qquad
 K(u)=C(u,1,Pu).
 \label{eq:eic-sign-countdown}
\end{equation}
Their defining equations imply
\begin{align}
 \operatorname{sg}(0)&=0,
 &\operatorname{sg}(Sx)&=1,\label{eq:eic-sign-computation}\\
 K(0)&=1,
 &K(Sx)&=x.\label{eq:eic-countdown-computation}
\end{align}

\paragraph{Constructor laws.}
Successor is injective: applying $P$ to $Sx=Sy$ gives $x=y$ by
\eqref{eq:eic-predecessor}.  Successor never equals zero.  Indeed,
$Sx=0$ gives $1=0$ by applying $\operatorname{sg}$ and using
\eqref{eq:eic-sign-computation}; symmetry and
\eqref{eq:eic-separation} give $\bot$.  Consequently $1\ne0$, where
$A\ne B$ abbreviates $A=B\to\bot$.  None of these arguments uses
induction.

\paragraph{The sign test.}
For every $u$,
\begin{align}
 &\operatorname{sg}(u)=0\lor\operatorname{sg}(u)=1,
       \label{eq:eic-sign-range}\\
 &\operatorname{sg}(u)=0\quad\longleftrightarrow\quad u=0,
       \label{eq:eic-sign-zero}\\
 &\operatorname{sg}(u)=1\quad\longleftrightarrow\quad u\ne0.
       \label{eq:eic-sign-one}
\end{align}
To prove the range assertion, split \eqref{eq:eic-coverage}.  The branch
$u=0$ gives sign $0$, and the branch $u=S(Pu)$ gives sign $1$.
The implication $u=0\to\operatorname{sg}(u)=0$ is computation.
For its converse, the successor branch of coverage would make
$\operatorname{sg}(u)=0$ imply $1=0$, a contradiction; in that branch
explosion gives $u=0$ as well.  Similarly, if
$\operatorname{sg}(u)=1$ and $u=0$, computation gives $0=1$.
Conversely, $u\ne0$ excludes the zero branch of coverage, and the
successor branch computes the sign as $1$.

Coverage and \eqref{eq:eic-selector} also give, for all $u,v$,
\begin{equation}
 C(u,v,v)=v.
 \label{eq:eic-selector-idempotence}
\end{equation}
Both coverage branches compute to $v$, so this identity again needs no
induction.

\paragraph{The countdown test and the notation $W$.}
For every $u$,
\begin{equation}
 K(u)=0\quad\longleftrightarrow\quad u=1.
 \label{eq:eic-countdown-zero}
\end{equation}
If $u=1=S0$, \eqref{eq:eic-countdown-computation} gives $K(u)=0$.
Conversely, assume $K(u)=0$ and split coverage.  If $u=0$, then
$K(u)=1$, contradicting separation, and explosion gives $u=1$.
If $u=S(Pu)$, computation gives $K(u)=Pu$; hence $Pu=0$ and
$u=S0=1$.

If a presentation uses a description named $W$ with defining equations
\[
 W(0)=1,\qquad W(Sx)=x,
\]
then it proves $K(u)=W(u)$ for every $u$ already from coverage.
In the zero branch both sides are $1$; in the branch $u=S(Pu)$ both
sides are $Pu$.  Thus $K$ and $W$ are provably equal on the whole domain,
even if they are distinct descriptions.  This equality does not appeal to
standard-model agreement or to uniqueness of an arbitrary recurrence.

\subsection{Four elementary uses of equation induction}

Choose descriptions for addition and truncated subtraction, with recursion
on the second argument:
\begin{align}
 x+0&=x, &x+Sy&=S(x+y),\label{eq:eic-addition}\\
 d(x,0)&=x, &d(x,Sy)&=P(d(x,y)).\label{eq:eic-subtraction}
\end{align}
Using these equations and \eqref{eq:eic-induction}, we now derive
\begin{align}
 0+y&=y,\label{eq:eic-left-zero}\\
 (Sx)+y&=S(x+y),\label{eq:eic-left-successor}\\
 d(Sx,Sy)&=d(x,y),\label{eq:eic-cancel-successors}\\
 d(Sa,a)&=1.\label{eq:eic-diagonal}
\end{align}
Each induction below has a single equation as its induction formula.

For \eqref{eq:eic-left-zero}, induct on $y$.  The base is $0+0=0$.
The step is $0+Sy=S(0+y)=Sy$.

For \eqref{eq:eic-left-successor}, hold $x$ fixed and induct on $y$.
At zero both sides are $Sx$.  At a successor the induction hypothesis
gives
\[
 (Sx)+Sy=S\bigl((Sx)+y\bigr)
       =S\bigl(S(x+y)\bigr)=S(x+Sy).
\]
Combining \eqref{eq:eic-left-successor} with the defining equation for
addition gives the useful rearrangement
\begin{equation}
 (Sy)+r=y+Sr.
 \label{eq:eic-shift-successor}
\end{equation}
Indeed, both sides are $S(y+r)$.

For \eqref{eq:eic-cancel-successors}, hold $x$ fixed and induct on $y$.
The base computes as
\[
 d(Sx,S0)=P(d(Sx,0))=P(Sx)=x=d(x,0).
\]
Applying $P$ to the induction equation gives the step:
\[
 d(Sx,S(Sy))=P(d(Sx,Sy))=P(d(x,y))=d(x,Sy).
\]

Finally, induct on $a$ for \eqref{eq:eic-diagonal}.  The base is
$d(S0,0)=S0=1$.  At a successor,
\[
 d(S(Sa),Sa)=d(Sa,a)=1,
\]
using \eqref{eq:eic-cancel-successors} and the induction hypothesis.

\subsection{An equation invariant that reflects actual equality}

The remaining arithmetic ingredient is the following equation, valid for
all $x,y$:
\begin{equation}
 C\bigl(\operatorname{sg}(d(x,y)),\ x,\ y+d(x,y)\bigr)=x.
 \label{eq:eic-residual-invariant}
\end{equation}
Its meaning is that a nonzero subtraction residual reconstructs the
original argument by addition.  The selector expresses this assertion as
one literal equation, so equation induction applies directly.

Hold $x$ fixed and induct on $y$.  At $y=0$, subtraction gives residual
$x$, and \eqref{eq:eic-left-zero} gives $0+x=x$.  The desired equation
is consequently $C(\operatorname{sg}(x),x,x)=x$, which is
\eqref{eq:eic-selector-idempotence}.

For the step, assume \eqref{eq:eic-residual-invariant} at $y$ and
abbreviate
\[
 r=d(x,y),\qquad r'=Pr=d(x,Sy).
\]
These are local term abbreviations, not additional function symbols.
Split coverage for $r'$.

If $r'=0$, then $\operatorname{sg}(r')=0$, and the successor instance
of \eqref{eq:eic-residual-invariant} reduces immediately to $x=x$.

In the other branch, $r'=S(Pr')$, so constructor separation gives
$r'\ne0$.  Split coverage for $r$.  Its zero branch would imply
$r'=P0=0$, a contradiction, from which the desired step equation follows.
In its successor branch,
\[
 r=S(Pr)=Sr',\qquad \operatorname{sg}(r)=1.
\]
The induction equation therefore computes to $y+r=x$.  Meanwhile,
\eqref{eq:eic-shift-successor} gives
\[
 (Sy)+r'=y+Sr'=y+r=x.
\]
Since $r'$ is a successor, $\operatorname{sg}(r')=1$, so this is
exactly the branch selected in the successor instance of
\eqref{eq:eic-residual-invariant}.  This completes the equation induction.
All case distinctions came from coverage; no induction on an implication
or a disjunction was used.

\paragraph{Reflection at residual $1$.}
We have, for arbitrary $a,b$,
\begin{equation}
 d(Sa,b)=1\quad\longleftrightarrow\quad a=b.
 \label{eq:eic-residual-reflection}
\end{equation}
For the forward direction, substitute $x=Sa$, $y=b$ into
\eqref{eq:eic-residual-invariant}.  If $d(Sa,b)=1$, the sign is $1$
and the selector returns its second branch, giving
\[
 b+1=Sa.
\]
The defining equation for addition gives $b+1=S(b+0)=Sb$.  Thus
$Sb=Sa$, and applying $P$ gives $b=a$, hence $a=b$.
Conversely, $a=b$ and \eqref{eq:eic-diagonal} give
$d(Sa,b)=d(Sa,a)=1$.

Define the equality characteristic by the fixed description
\begin{equation}
 e(a,b)=\operatorname{sg}\bigl(K(d(Sa,b))\bigr).
 \label{eq:eic-equality-characteristic}
\end{equation}
The sign, countdown, and residual reflection lemmas give
\begin{align}
 e(a,b)=0&\quad\longleftrightarrow\quad a=b,
       \label{eq:eic-equality-reflection}\\
 e(a,b)=0&\ \lor\ e(a,b)=1.
       \label{eq:eic-equality-range}
\end{align}
In detail, $e(a,b)=0$ is equivalent to $K(d(Sa,b))=0$ by
\eqref{eq:eic-sign-zero}, which is equivalent to $d(Sa,b)=1$ by
\eqref{eq:eic-countdown-zero}, and then to $a=b$ by
\eqref{eq:eic-residual-reflection}.  The range assertion is
\eqref{eq:eic-sign-range} applied to $K(d(Sa,b))$.

The reflection property concerns the logic's actual equality throughout
the domain.  It was not assumed as part of a numerical comparison symbol,
and its proof did not restrict $a$ or $b$ to standard numerals.
It also yields decidable equality.  In the $e(a,b)=0$ branch,
\eqref{eq:eic-equality-reflection} gives $a=b$.  In the $e(a,b)=1$
branch, an assumption $a=b$ would give $e(a,b)=0$ and therefore $0=1$,
contradicting separation.  Hence
\begin{equation}
 \forall a\,\forall b\,\bigl(a=b\lor(a=b\to\bot)\bigr)
 \label{eq:eic-equality-decidable}
\end{equation}
is a theorem, not an extra hypothesis.

\subsection{Compilation of all quantifier-free formulas}

For each fixed quantifier-free formula $A(\vec p)$, construct a term
$b_A(\vec p)$ by external recursion on its finite syntax.  Take equality,
$\bot$, $\land$, $\lor$, and $\to$ as the primitive formula constructors;
negation is implication into $\bot$.  Set
\begin{align}
 b_{s=t}&=e(s,t), & b_{\bot}&=1,\label{eq:eic-boolean-atoms}\\
 b_{A\land D}&=C(b_A,b_D,1),\label{eq:eic-boolean-and}\\
 b_{A\lor D}&=C(b_A,0,b_D),\label{eq:eic-boolean-or}\\
 b_{A\to D}&=C(b_A,b_D,0).\label{eq:eic-boolean-implies}
\end{align}
Every term here is a finite composition of fixed primitive-recursive
descriptions; if desired, projections and composition give it a single
description symbol.  No internal quantification over formulas or programs
is involved.

An external induction on the syntax of $A$ proves, within the arithmetic
theory,
\begin{equation}
 \bigl(b_A=0\lor b_A=1\bigr)
 \quad\text{and}\quad
 \bigl(b_A=0\longleftrightarrow A\bigr).
 \label{eq:eic-boolean-correctness}
\end{equation}
Here are the full connective checks.  In every case $b_A=1$ implies
$\neg A$: if $A$ held, the already established implication
$A\to b_A=0$ would contradict separation.

\paragraph{Atomic formulas and falsity.}
For an equation, both claims are
\eqref{eq:eic-equality-reflection}--\eqref{eq:eic-equality-range}.
For falsity the range assertion is $1=0\lor1=1$, with its right disjunct
proved by reflexivity.  Also $1=0\to\bot$ follows from separation and
symmetry, while $\bot\to1=0$ follows by explosion.  Thus $1=0$ is
equivalent to $\bot$, as required.

\paragraph{Conjunction.}
Split $b_A=0\lor b_A=1$.  In the zero branch, $A$ holds and
$b_{A\land D}=b_D$.  Its range follows from the range of $b_D$,
and its being zero is equivalent to $D$, hence to $A\land D$ under
the known $A$.  In the one branch, $\neg A$ holds and
$b_{A\land D}=1$.  Its range is immediate.  If it were zero,
separation would give $\bot$, hence $A\land D$; conversely an
assumption $A\land D$ contradicts $\neg A$, and explosion gives
$b_{A\land D}=0$.

\paragraph{Disjunction.}
In the branch $b_A=0$, the formula $A$ holds and
$b_{A\lor D}=0$.  The range is immediate, and both $b_{A\lor D}=0$
and $A\lor D$ hold, proving their equivalence.  In the branch
$b_A=1$, we have $\neg A$ and $b_{A\lor D}=b_D$.  Its range follows
from the range of $b_D$.  If its value is zero, then $D$ and hence
$A\lor D$.  Conversely, from $A\lor D$, the $A$ alternative
contradicts $\neg A$ and yields $b_D=0$ by explosion, while the
$D$ alternative yields $b_D=0$ directly.  This proves the required
equivalence in both alternatives.

\paragraph{Implication.}
In the branch $b_A=0$, the formula $A$ holds and
$b_{A\to D}=b_D$.  Its range follows from the range of $b_D$.
If its value is zero, then $D$, which gives $A\to D$.  Conversely,
$A\to D$ and the known $A$ give $D$, hence $b_D=0$.
In the branch $b_A=1$, we have $\neg A$ and
$b_{A\to D}=0$.  Its range is immediate.  To prove $A\to D$,
assume $A$; this contradicts $\neg A$, and explosion gives $D$.
Thus both $b_{A\to D}=0$ and $A\to D$ hold in this branch, and
their equivalence follows.  This argument uses the proved two-value
range of the finite characteristic.  It does not use excluded middle
for arbitrary formulas or change the surrounding intuitionistic logic.

This completes the external syntactic induction proving
\eqref{eq:eic-boolean-correctness} for every quantifier-free formula.

\subsection{Recovery of induction and the full theory}

Let $A(n,\vec p)$ be any fixed quantifier-free formula.  Suppose an
arbitrary retained context $\Gamma$ proves its base and its induction
step:
\[
 \Gamma\vdash A(0,\vec p),
 \qquad
 \Gamma,A(z,\vec p)\vdash A(Sz,\vec p),
\]
with the usual freshness condition on $z$.  By
\eqref{eq:eic-boolean-correctness}, the base gives
$\Gamma\vdash b_A(0,\vec p)=0$.  Under the additional assumption
$b_A(z,\vec p)=0$, the same equivalence gives $A(z,\vec p)$,
the given step gives $A(Sz,\vec p)$, and the equivalence gives
$b_A(Sz,\vec p)=0$.  Apply literal-equation induction to
$b_A(n,\vec p)=0$ and then apply its reflection implication.
The result is $\Gamma\vdash A(t,\vec p)$ for every target term $t$.
Thus the theory proves the full contextual quantifier-free induction
rule, including in arbitrary quantified contexts.

\paragraph{Completion theorem.}
In the common language of fixed primitive-recursive descriptions, the
theory
\[
 \mathsf{PRE}+\mathsf{Cov}+\mathsf{Sep}+\mathsf{EI}
\]
has exactly the same sequents as ordinary first-order intuitionistic
primitive recursive arithmetic with all quantifier-free induction
instances.  The forward simulation follows from the recovered
quantifier-free induction, successor injectivity, and successor separation;
all primitive-recursive defining equations and all logical rules are
shared.  For the reverse simulation, literal-equation induction is an
instance of quantifier-free induction, and successor separation gives
$\mathsf{Sep}$.  Quantifier-free induction proves $\mathsf{Cov}$:
at zero choose its left disjunct, and at a successor use
$Sx=S(P(Sx))$, which follows from $P(Sx)=x$, for its right disjunct.

Both interpretations therefore use the full domain, defined by $x=x$,
and preserve every function symbol, equality, logical connective, and
quantifier.  Both composites have the identity isomorphism, with graph
$x=y$.  The theories are consequently identity-domain bi-interpretable
and equiconsistent.  This completion theorem is unconditional once the
literal-equation induction rule has been obtained; it has no dependency
on a classical successor-uniqueness reduction.

\section{Pairing, primitive recursion, and internal numerical composites}
\label{sec:iterator-coding}
This appendix is adapted with attribution from the progress-closure
arithmetic development~\cite{progress}. Its prerequisites have now been
proved: scalar predecessor, arithmetic reflection, and full quantifier-free
induction. Pairing is a derived proof tool; it is not part of any checker.

\subsection{Pairing from a bounded square-root search}

Once QFI is available, a familiar square-based injection gives a
particularly direct compiler. Write $a\le b$ for $d(a,b)=0$.
The usual linear-order and semiring identities follow by fixed
quantifier-free inductions from the successor, addition, and monus
equations. Multiplication itself is the scalar term
\[
                        a\cdot b=I(b,0;u.u+a).
\]
Only primitive-recursive arithmetic is involved in the following
bounded root search; there is no minimization over an unbounded domain.

For a fixed parameter $n$, put
\[
 U_n(v)=C(d((Sv)^2,n),Sv,v),\qquad
 R_n(i)=I(i,0;v.U_n(v)),\qquad q(n)=R_n(n).
\]
The update increments $v$ exactly when $(v+1)^2\le n$, and
otherwise retains it. QFI in $i$, with $n$ fixed, proves
\[
 R_n(i)\le i,\qquad R_n(i)^2\le n,\qquad
 R_n(i)=i\ \vee\ n<(S(R_n(i)))^2.              \tag{Root invariant}
\]
The base has $R_n(0)=0$. If the update increments, its test
contradicts the strict-upper-bound alternative in the third invariant.
Thus the old state equals $i$, and the new state equals $Si$;
its square bound is the increment test. If the update does not increment, the
strict upper square bound holds and the state remains fixed.
An already established strict upper square bound forces every
subsequent update to remain fixed. These are exactly the cases
needed for the three displayed invariants.

At $i=n$, either the strict upper square bound already holds or
$R_n(n)=n$, in which case $n<(n+1)^2$ supplies it. Thus
\[
                        q(n)^2\le n<(q(n)+1)^2.           \tag{Root}
\]
Monotonicity of squares shows uniqueness of a number satisfying
Root. All formulas just used are quantifier free in the scalar
language, with fixed $n$ during the induction.

Now define
\[
 \pi(a,b)=(a+b)^2+a,\qquad
 L(n)=d(n,q(n)^2),\qquad H(n)=d(q(n),L(n)).
\]
For $k=a+b$, one has $k^2\le k^2+a<(k+1)^2$ since $a\le k$.
Root uniqueness gives $q(\pi(a,b))=a+b$, and therefore
\[
                    L(\pi(a,b))=a,\qquad H(\pi(a,b))=b.
\]
Surjectivity of $\pi$ is unnecessary. The original numerical domain
is never replaced by the range of this coding; pairing is used only
to store the finite state of a fixed primitive recursion. This
standard square pairing is a convenient reconstruction here, not
a claimed new arithmetic basis.

\subsection{Compiling each primitive recursion}

Suppose the fixed target description is
\[
 f(0,\vec p)=g(\vec p),\qquad
 f(Sn,\vec p)=h(n,f(n,\vec p),\vec p).
\]
Assume by external structural induction that the earlier $g,h$
already have scalar translations $g^*,h^*$. Define
\[
 W(u,\vec p)=\pi(SL(u),h^*(L(u),H(u),\vec p)),
\]
\[
 V(n,\vec p)=I(n,\pi(0,g^*(\vec p));u.W(u,\vec p)),
 \qquad f^*(n,\vec p)=H(V(n,\vec p)).
\]
The update has only its state variable bound, with every $p_i$
fixed. Equation induction proves $L(V(n,\vec p))=n$.
The two projection equations and scalar computation then prove
both defining equations for $f^*$ internally. Initial functions
and composition translate directly. This completes the translation
of every fixed finite PR description without adding a universal
evaluator.

\subsection{Why both numerical round trips are internal}

Two fixed scalar iterators with the same seed and provably equal
updates are internally equal at all counters: use literal-equation
induction comparing their values. Their base agrees, and their
successor equations plus congruence prove the step. The same
uniqueness argument works for ordinary primitive recursion.

For an ordinary PR description, induct externally on its finite
construction. The double translation satisfies its original
defining equations, so internal equation induction identifies it
with the original operation. For a native scalar description,
earlier translated updates agree by the external induction
hypothesis; scalar uniqueness identifies the double-translated
iterator. These proofs hold at arbitrary internal counters.

Thus the domain relation $x=y$ is provably a total bijection
preserving every fixed numerical operation in both composites.
The formula constructions and round trips in
Section~\ref{sec:choice-introduction} give the corresponding formula
composites. Both interpretations preserve falsity, taking native
$0=1$ to an ordinary contradiction and ordinary $\bot$ to native
$0=1$. Finite proof translations therefore establish equiconsistency.

\section{A uniform arity encoding and its checking theorem}
\label{sec:arity-checker}

The arithmetic and object logic remain the scalar-iteration Eq/$\D$
calculus specified above.  The change is in its finite certificate
representation.  A paired formula is unary, an equation is binary, and
an internal closure is ternary:
\[
 \D_{x,y}B\longleftrightarrow[B],\qquad
 s=t\longleftrightarrow[s,t],\qquad
 \Q_yB\longleftrightarrow[B,B,B].
 \tag{Arity encoding}\label{eq:arity-encoding}
\]
The three copies in a closure are representation redundancy: only the
third has mathematical content.  Their benefit is that tripling a unary
major constructs exactly its closure.  Tripling any other nonempty
formula array produces an array of arity at least six, which will denote
truth.  This makes the other major cases neutral without a separate
test.

The previous generated-assumption and constructed-computation interfaces
are inherited from the corpus~\cite{generated,constructed}.  The new
encoding makes a single arithmetic expression, $5-\operatorname{arity}$,
perform the necessary depth adjustments.  Its soundness depends on an
explicit interpretation of surplus syntax.  It would not be sound to
drop arity checks while continuing to treat every surplus array as an
unconstrained new connective.

\subsection{The complete sorted interpretation of arrays}
\label{subsec:sorted-arrays}

Write $\top$ for the fixed equation $0=0$.  A protected numerical slot
has width two.  The numerical root has depth three: the coordinates
$0,1,2$ all denote zero, and under $p$ protected slots the depth is
$3+2p$.  Coordinates $3+2j$ and $4+2j$ name slot $j$, counted from
the innermost slot, for $0\leq j<p$.  Formula depth is $2p$ under the
same $p$ slots.  Negative coordinates admitted by the current global
floor name rigid numerical parameters.

The decoder is sorted by the position in which a tree is used.  Let
$\rho$ be its ordered environment of protected numerical slots.
The notation $(a,\rho)$ means prepend a slot of value $a$.
A superscript $\sharp$ denotes a child decoded in the indicated
environment.

\begin{center}
\begin{tabular}{cll}
\toprule
Numerical arity & Meaning & Relevant child environment\\
\midrule
$0$ & $0$ & No child\\
$1$: $[b]$ & $Sb^\sharp$ & $(0,0,\rho)$\\
$2$ & $0$ & Annotations ignored\\
$3$: $[n,a,h]$ & $I(n^\sharp,a^\sharp;v.h^\sharp)$ & $(0,\rho)$ for $n,a$; $(v,\rho)$ for $h$\\
$\geq4$ & $0$ & All children ignored\\
\bottomrule
\end{tabular}
\end{center}

\begin{center}
\begin{tabular}{cll}
\toprule
Formula arity & Meaning & Relevant child environment\\
\midrule
$0$ & $\top$ & No child\\
$1$: $[B]$ & $\D_{x,y}B^\sharp$ & $(y,x,\rho)$\\
$2$: $[s,t]$ & $s^\sharp=t^\sharp$ & $\rho$ numerically\\
$3$: $[P,U,B]$ & $\Q_yB^\sharp$ & $(y,\rho)$ for $B$\\
$\geq4$ & $\top$ & All children ignored\\
\bottomrule
\end{tabular}
\end{center}

The first two entries of a ternary formula are annotations.  A canonical
closure copies its body into both of them, as in
\eqref{eq:arity-encoding}.  Successfully admitted scalar formula padding
also denotes $\top$.  An internal closure has the ordinary meaning
$B(0)\to\forall yB(y)$; it is not an equation and is not eligible for
compound induction.

Canonical encoded object expressions use arrays and integer coordinates,
respect these slots, and contain no high-arity node.  The checker itself
still generates high-arity truth entries in administrative contexts,
even during a canonical proof.  In particular, constructing a
successor beneath positive bound coordinates requires passing through
two dummy zero slots.  Constructing an iterator requires one new slot:
zero in its counter and seed, and its genuine state variable in the
update.  Rigid negative parameter names are unaffected by these shifts.

The additional raw shapes are certificate aliases with the displayed
meanings.  They are not independent predicate or function symbols added
to the object theory.  Our exact completeness statement concerns the
canonical native language.  We do not require every possible raw spelling
of a semantic alias to be interchangeable under the same raw certificate.

\subsection{One depth calculation replaces the lookup table}

For every array the opener recursively adds $5-\operatorname{arity}$
to its depth.  The meaningful cases are
\[
\begin{array}{c|ccc}
 \operatorname{arity}&1&2&3\\\hline
 \text{depth increment}&4&3&2.
\end{array}
\]
Unary $\D$ adds two numerical binders; unary successor adds two dummy
zero slots.  Binary equality changes formula depth $2p$ to numerical
depth $3+2p$.  Ternary $\Q$ adds one numerical binder; ternary iteration
adds one state slot, or one dummy slot in its counter and initializer.
The same calculation therefore covers all relevant constructors.

Binary numerical annotations denote zero, so their parity-changing
traversal has no mathematical effect.  Empty arrays have no children.
At arity at least four, the entire node already denotes a constant.
Even a zero or negative depth increment below such a node cannot expose
one of its ignored children as a meaningful variable.  Traversal may
still reject an alias; it cannot change its stipulated interpretation.

This is a semantic assertion, not raw idempotence of the opener.
For example, at floor $-1$, opening a six-element array of zeros with
image $-1$ can yield six copies of $-1$, while a later no-image opening
of that result fails at a negative child depth.  Both formula arrays
denote $\top$.  The proof below only concerns successful calls at meaningful
sorted positions.  Canonical input object expressions contain no surplus
arrays; the administrative exports that generate them are justified
separately as truth entries.

\subsection{Opening and scope}
\label{subsec:opening-lemma}

Write $i(a,t,d)$ for the executable opener at floor $k$.  It preserves a
scalar when $k\leq a<d$.  At a numerical position, every other relevant
scalar is the distinguished substitution hole; the supplied image is
checked at numerical root depth three before insertion.  At a formula
position scalar padding is truth and carries no numerical hole.
Missing images at used numerical holes reject.

\begin{lemma}[Semantic opening]
\label{lem:semantic-opening}
For every successful call at a meaningful formula root of depth $2p$
or numerical root of depth $3+2p$, with $p\geq0$, decoding its result
is capture-free filling of the distinguished hole in the decoded input
template.  These depths protect $p$ numerical slots; admitted rigid
parameters are preserved.  If an image is semantically unused, zero
may be chosen as its mathematical representative.  Recursive calls
at negative depths inside ignored annotations need no independent
sorted interpretation.
\end{lemma}

\begin{proof}
Induct on the finite raw tree.  At a relevant scalar, the interval
comparison is exactly the protected-coordinate test.  At a meaningful
unary, binary, or ternary node, the preceding tables supply precisely
the required additional slots.  At numerical binary annotations,
high-arity nodes, and ignored closure fields, the relevant mathematical
meaning is constant or absent, so the conclusion is immediate
regardless of traversal inside the annotation.

An inserted image is checked at its own numerical root.  Its free
parameters are rigid negative names, and it has no unbound positive
coordinate there.  Insertion beneath protected slots therefore captures
none of its parameters; its internal binders retain their relative
indices.  This proves the inductive step for substitution as well as
for mere traversal.
\end{proof}

A successful no-image check consequently preserves the interpretation
of its checked closed expression.  Raw structural equality of checked
outputs implies equality of their decoded meanings.  The checker does
not infer equality merely from a semantic comparison of two different
programs or numerical terms.

The global floor decreases by one at each non-hypothesis continuation.
If its old value is $K$, the integer $K-1$ is the designated fresh
eigenparameter.  With integer coordinates it is the only new name.
If fractional coordinates are accepted, other names in $[K-1,K)$ may
also occur in the finite certificate.  These are additional fresh
parameters.  They must not be silently equated to the designated
integer eigenparameter.

Previously checked relevant goals and retained hypotheses mention none
of these new names: their earlier no-image checks would reject a
relevant numerical occurrence below the old floor.  Coordinates inside
ignored annotations do not contribute free variables to the decoded
formula.  This is the freshness invariant used in every rule below.

\subsection{A computation hypothesis only needs a zero instance}
\label{subsec:zero-supply}

The shared continuation can open a generated arithmetic equation at a
fresh parameter.  That equation need not be a theorem uniformly in the
new parameter.  The following admissible rule is sufficient.

\begin{lemma}[Supply with a known zero instance]
\label{lem:zero-supply}
Suppose $\vec a$ is fresh for $\Gamma,F$, the calculus proves
$H(\vec0)$, and $\Gamma,H(\vec a)\vdash F$.  Then
$\Gamma\vdash F$.
\end{lemma}

\begin{proof}
Substitute zero for every member of $\vec a$ in the second derivation.
Freshness leaves $\Gamma,F$ unchanged.  Cut in the derivation of
$H(\vec0)$.  Both substitution and cut are admissible in the native
calculus.  This is a finite proof transformation and does not use a
consistency assumption.
\end{proof}

All names newly admitted when the floor changes belong to $\vec a$;
for noninteger aliases this includes the finite collection in
$[K-1,K)$ that actually occurs.  On decoding, their simultaneous zero
substitution agrees with opening the generated expression at zero using
the old floor.  Previously admitted parameters are preserved.  New
rigid names and unprotected placeholders both become zero.  The opener
may keep different raw zero spellings, but their decoded values agree.

The constructed-computation packet supplies a triple $s=[n,b,r]$ and a
mode $u$.  In successor mode its generated equation has raw form
\[
 \bigl[\,[[n],b,r],\ i(r,s,3)\,\bigr],
 \tag{Step supply}\label{eq:step-supply}
\]
while in zero mode it has form
\[
                    \bigl[\,[u,b,r],\ b\,\bigr].
 \tag{Zero supply}\label{eq:base-supply}
\]
Here successor mode means that $u$ is truthy; the zero mode has falsy
$u$.  These expressions describe actual arrays, before their final
opening.  In \eqref{eq:step-supply} the counter is constructed as a
successor, rather than accepted on an unsupported assertion that it is
one.  This construction is inherited from the corresponding draft
interface~\cite{constructed}.

\begin{lemma}[Sound constructed computation]
\label{lem:constructed-supply}
Every successful arithmetic supply has a decoded zero instance
derivable from the two scalar-iteration computation schemas.
\end{lemma}

\begin{proof}
First apply the simultaneous zero substitution described above.  In
zero mode every successful falsy numerical counter denotes zero, so
\eqref{eq:base-supply} is an $I0$ instance.  In the exact-list profile
the relevant accepted aliases are an empty array, numeric zero, false,
and numerically equal zero values.  Null, strings, and dictionaries
cannot enter a meaningful numerical position through an unchecked
coercion.

In successor mode the left counter is an actual successor.  If the
update uses its state in a semantically relevant position, opening it
invokes the no-image check of the recursive source $s$ at numerical
root three.  This checks the entire source, including its update.
Only coordinates three and four are aliases for its single state slot.
A genuinely escaping coordinate in that update is revisited during
the recursive-image check and rejects.  Its counter and initializer
are likewise closed with their prescribed dummy zero slot.

Inserting the extra dummy slots of the displayed successor does not
alter such a checked predecessor: extra unbound slots have value zero,
and actual free parameters are rigid.  Lemma~\ref{lem:semantic-opening}
therefore identifies the right side of \eqref{eq:step-supply} with the
update applied to exactly the recursive predecessor value occurring
on its left.  This is the $IS$ schema.

If the update has no semantically relevant state occurrence, its
decoded value is constant in that variable.  Its zero-normalized value
is the same in the source update and on the right, independently of
the unused recursive image; this is again an $IS$ instance.  Traversal
inside an ignored annotation may still inspect the image and reject.
That removes a certificate and does not invalidate the argument for a
successful one.

Finally, successful three-field source extraction cannot hide another
JSON operator.  A scalar is not iterable.  A three-character string,
or a three-key JSON dictionary, necessarily puts a nonempty string in
a relevant source field, which the inspected opener rejects.  The
successful source is therefore the stated array triple.
\end{proof}

The recursive-image check in the state-dependent case is essential.
Merely opening a final displayed equation at zero would not establish
that its two state interpretations agree.  Combining
Lemma~\ref{lem:constructed-supply} with Lemma~\ref{lem:zero-supply}
justifies both actual image choices of the common continuation, including
the branch that retains a fresh-parameter hypothesis.

\subsection{Audit of the logical continuations}
\label{subsec:logical-audit}

\paragraph{Hypothesis closure.}
The initial opener check and membership comparison select the same
decoded context formula.  Additional aliases can make a spelling less
convenient, but cannot identify an equation with a false denotation
solely by semantic evaluation.

\paragraph{Equality transport.}
The major must match $[s,u]$, and opening the proposed formula template
at $u$ must match the target.  The continuation checks its opening at
$s$.  Lemma~\ref{lem:semantic-opening} gives ordinary transport along
the decoded equality.  The argument applies to compound formulas,
internal clauses, and truth aliases as well as equations.

\paragraph{Closure specialization.}
The closure pattern is the triplicate $[t,t,t]$.  The guard forces the
old instance to have numerical image zero; the new instance may have
an arbitrary numerical image.  This is exactly
$B(0)\to\forall yB(y)$.  The zero guard is indispensable.  Without
it, a closure with false zero instance and true successor instances
would turn a proved successor instance into falsity.

\paragraph{Paired elimination.}
The export is always three concatenated copies of the major.  A unary
major $[B]$ therefore exports $[B,B,B]$.  All three copies are opened
identically; the third copy has precisely the witness instance required
by paired elimination.  The newly allocated witness is fresh for the
retained context and conclusion.

An equation major instead exports arity six; a closure major exports
arity nine; a larger array exports still larger arity.  These are all
truth aliases.  Thus failed elimination endpoint claims have only the
effect of adjoining a theorem.  A numeric scalar major, if accepted as
formula padding, exports scalar padding after multiplication and also
denotes truth.  Unsupported types reject.

A ternary formula with unequal annotations still denotes the closure
of its third child.  It fails the triplicate specialization pattern
unless the three fields match.  This loses a possible raw spelling,
not a mathematical soundness condition.  Canonical compilation uses
triplicate closures throughout.

\paragraph{Paired introduction.}
The call \texttt{i(*a,t,2)} uses the function's argument count to enforce
the relevant outer arity.  For a unary goal $a=[B]$, it is exactly
\texttt{i(B,t,2)}: the universal slot is protected while the
existential witness is opened.  For arity at least two, there are too
many arguments and the call rejects.  An empty goal can instead make
the legal two-argument call \texttt{i(t,2)}.  Its original target was
already $\top$, so accepting it cannot add a false theorem.

For the unary goal, the subsequent successor opening gives the intended
premise $B(t,Sz)$.  The local hypothesis is $B(t,0)$ in case mode and
$B(t,z)$ in progress mode.  The freshness invariant supplies the
usual eigenvariable conditions.

\paragraph{Induction remains restricted.}
Among meaningful nonconstant matrix constructors, only binary equality
has even arity.  Unary paired formulas and ternary closures have odd
arity and receive case introduction.  Empty or surplus arrays that
reach the progress branch denote the fixed equation $0=0$, whose
induction is valid.  Scalar padding that cannot provide a length
rejects.  The broader certificate alphabet therefore adds no induction
on a quantified matrix.

\paragraph{Additional fresh parameters.}
An accepted subderivation can contain newly admitted fractional names,
but the older context and target do not depend on them.  Replace these
auxiliary names by zero when closing the rule, retaining its designated
integer eigenvariable where that rule needs one.  The common opening
of an already checked retained goal leaves its decoded meaning
unchanged.  This handles the actual global-floor convention without
pretending that every admissible negative number is an integer.

\subsection{Canonical completeness, including zero witnesses}
\label{subsec:canonical-completeness}

Rename a finite native proof's eigenvariables to fresh negative integers.
Encode its terms with the slot conventions of
Section~\ref{subsec:sorted-arrays}; encode paired formulas as unary and
closures as triplicate.  Hypothesis leaves consume no instruction.
Transport, closure specialization, and paired elimination have their
displayed packets.  Arithmetic computation leaves use
\eqref{eq:base-supply} or \eqref{eq:step-supply} and then close the
target by hypothesis lookup.

One representational issue requires an actual proof transformation.
The packet's truthy witness field selects introduction.  A literal zero
witness therefore cannot simply be inserted into that field.  Define
\[
                 J(t)=I(0,t;v.0).
\]
The computation axiom proves $J(t)=t$, and $J(t)$ has a nonempty array
representation even when its value is zero.  To introduce with a falsy
witness spelling $t$, introduce instead with $J(t)$: transport its
retained assumption to the assumption using $t$, use the original
premise derivation, and transport the resulting successor instance
back to the instance using $J(t)$.  Derived symmetry supplies the
reverse transport.  This is a finite transformation of the premise,
and the paired conclusion itself is unchanged.  If the existential
coordinate is unused, any closed truthy witness works without a
conversion.

The scalar coordinate $1$ is also a semantic zero at numerical root
three, but that observation alone is not a completeness proof.  Closure
specialization and case introduction insert literal raw zero, and
distinct zero aliases need not be interchangeable under a fixed raw
proof spelling.  The computation-backed box gives the required
derivation instead of assuming representational interchangeability.

The mode zero selects progress at an equation.  A different falsy
mode, such as the empty array, selects case.  These are certificate
control tags, not numerical assertions about the witness.  All steps
in the compilation are finite and effective.

\begin{theorem}[Conservativity and completeness of the arity checker]
\label{thm:arity-checker}
Every accepted certificate has its decoded conclusion provable in the
ordinary realization of the native calculus.  For every canonical
native end sequent, the internal checking judgment and the native
calculus have the same consequence relation, in the idealized
unbounded integer-and-array implementation.  In particular, the public
empty-context driver recognizes exactly the canonical closed theorems.
The translation in each direction is effective on finite derivations.
\end{theorem}

\begin{proof}
For soundness, induct on the finite accepting call tree.  Hypothesis
closure uses Lemma~\ref{lem:semantic-opening}; the logical cases are
the audited continuations above.  Arithmetic supply uses
Lemmas~\ref{lem:constructed-supply} and \ref{lem:zero-supply}.
Every discarded annotation is accounted for by its zero or truth
interpretation, and every fresh parameter by the scope invariant.
For completeness, compile a finite native derivation by the preceding
construction, including the boxed-witness transformation.  Its
canonical input expressions require no surplus-arity aliases or fractional
names; generated administrative surplus arrays have exactly the truth
meaning already proved.  This proves both implications without using a finite test
suite as a substitute for induction on derivations.
\end{proof}

The unchanged native calculus is bi-interpretable with full first-order
$\ipra$ by the inherited arithmetic and logical bridge
\cite{generated,corpus}.  In particular,
\[
 (\D_{x,y}B)^\circ
 =\exists x\bigl(B^\circ(x,0)\to\forall yB^\circ(x,y)\bigr).
\]
The whole numerical domain and actual equality are retained.  Its
internally proved numerical round trips and its formula-composite
derivations therefore transfer through
Theorem~\ref{thm:arity-checker}.  We have changed the checked
representation, not replaced bi-interpretability by standard-model
agreement or supplied a new proof of the scalar predecessor theorem.

\subsection{JSON and finite-runtime qualifications}

The primary profile uses an ordinary list-type test.  The smaller byte
profile instead tests multiplication by zero against an empty list.
Among standard JSON runtime values the latter recognizes exactly lists:
numeric values and booleans produce numeric zero, strings produce an
empty string, and null and dictionaries raise rather than becoming
arrays.  This argument assumes ordinary decoded JSON values, not
arbitrary Python objects with overloaded methods.

Finite nonnegative real coordinates have the same interval
interpretation: $[0,3)$ is zero, and $[3+2j,5+2j)$ names slot $j$.
Finite negative coordinates in the admitted range are additional rigid
names.  Equal Python numerical values have the same decoded role,
including Boolean aliases and signed zero.  Nonfinite values admitted
by Python's permissive JSON reader cannot survive a relevant closed
numerical check.  They may be filled as template holes or ignored
inside a constant annotation; neither makes them unchecked parameters.

The alternative opener based on generic iteration has a separate
runtime qualification.  Empty strings and empty dictionaries traverse
as empty arrays; a dictionary whose sole key is empty traverses as a
unary array with an empty child.  Dictionary values are not visited.
Nonempty strings can rely on CPython's recursion limit to reject, as
also occurs in a current corpus comparison profile~\cite{transient}.
The idealized structural-termination claim for the exact-list profile
must not be transferred to that generic profile unchanged.

All executable profiles have ordinary finite memory, stack, and input
limits.  These can reject a sufficiently large otherwise valid proof.
The formal consequence relation is the unbounded finite-certificate
calculus; the delivered program checks certificates within its runtime
resources.  The measurements charge the entire delivered checker,
while the semantic proof and canonical producer are supplied separately.

% Independent audit fragment for the corrected array-truth interface.
% Reviewed files: kernel_notes/array_truth_checker{,_bytes,_joint}.py.
% IMPORTANT: the final mode comparison is with u, not s.
% Requires only ordinary LaTeX plus amsmath/amssymb.

\section{Why the absent major can be a proved empty array}
\label{sec:neutral-major-audit}

This section audits the new control-flow change on raw input, including
noncanonical JSON.  It uses the inherited sorted arity decoder and
semantic-opening lemmas for the numerical and logical rules.  It does not
replace those lemmas with finite testing or claim proof-assistant
certification.

\subsection{The exact change}

The initial context is the one-element list \texttt{[[]]}.  Its sole
formula, the empty array, denotes $0=0$ under the inherited decoder.
At every non-hypothesis command the checker now executes
\begin{verbatim}
    m, s, u, t = q()
    c(m, g)
\end{verbatim}
without testing whether \texttt{m} is nonempty.  A canonical command with
no substantive major puts \texttt{[]} in this field.  Its premise closes
immediately by hypothesis, without consuming a command or changing the
fresh-name floor.

The final choice of the retained assumption uses
\begin{verbatim}
    (m or len(a) % 2 == u) and k
\end{verbatim}
instead of comparing the arity parity with \texttt{m}.  In an
introduction command, \texttt{u=0} selects equation progress and
\texttt{u=[]} selects the zero-case premise.  This reuse is legitimate
because introduction is reached only when both \texttt{m} and
\texttt{u} are false-valued.  Computation still requires a truthy
initializer \texttt{u}; its missing literal-zero cases have the
inherited certified-zero repair.

The source is deliberately left with the old matcher
\begin{verbatim}
    m in ([s, u], not s and [t] * 3)
\end{verbatim}
and no added \texttt{m and} guard.  As an isolated Python expression this
would be unsafe: if \texttt{s} is truthy, its second alternative is
\texttt{False}, which compares equal to numeric zero.  The following
whole-program invariant proves that this dangerous major cannot reach
the matcher.

\subsection{Two facts about the root opener}

Call a Python value \emph{numeric} here when it is an integer, a float,
or a Boolean returned by the JSON loader.  The finite Python execution
also admits the loader's optional nonfinite float values.  No custom
numeric class is supplied by a certificate.

Write $O_k(r;t)$ for the opener on the numeric value $r$, with image $t$
and its default formula-root depth zero.  Direct evaluation of the
program gives
\[
 O_k(r;t)=
 \begin{cases}
 r,&k\le r<0,\\
 0,&\text{otherwise}.
 \end{cases}
 \tag{Root scalar}
\]
In particular this result is numeric and independent of the image $t$.
The formula covers negative allocated names and negative fractional
names without identifying different numerical values.  Signed zero and
Booleans have Python's ordinary numerical equality behavior.  A NaN
fails the interval test and returns zero; the relevant infinite values
also return zero.  None of these results is a list.

A second fact is that a successful root opening of a list returns a
list.  Under either strict list predicate, no other input type can
return a numeric result except a numeric input.  Under the generic
iterable opener, a successful iterable traversal also returns a list;
if traversal fails, the numerical comparison cannot succeed on a list,
string, dictionary, or null.  Thus in every profile:
\begin{quote}
If a root opening returns a numeric result, its argument was numeric;
if it succeeds on a list, its result is a list.
\end{quote}
These statements concern finite successful evaluations.  They assert
neither arbitrary revalidation idempotence at ignored annotations nor
unlimited-stack termination for the generic iterable profile.

\subsection{No numeric goal can be proved from list-valued assumptions}

\paragraph{Lemma.}
For any initial integer fresh-name floor and any finite remaining
certificate, there is no finite successful invocation
\texttt{c(a,g)} whose raw goal \texttt{a} is numeric and whose every
context element is a list.

\emph{Proof.}
Suppose otherwise, and select such a successful invocation with the
fewest calls to \texttt{c} in its completed recursive trace.  Its first
hypothesis check cannot close: the root opening of its numeric goal is
numeric, and no numerical JSON value equals a Python list.

The invocation must therefore consume a command and successfully return
from \texttt{c(m,g)}.  If \texttt{m} were numeric, this would be a
strictly shorter successful invocation of the forbidden kind.  Thus the
major is not numeric.

In either strict profile, a successfully checked nonnumeric major is a
list.  In the generic profile an empty string or an empty dictionary can
also open to the empty list.  A nonempty dictionary can occasionally
have its iterable alias proved, but it cannot occur as a truthy major
in the selected completed parent invocation: the subsequent expression
\texttt{m * 3} raises a type error.  A nonempty string has no finite
successful opening.  Consequently every truthy major that survives to
the continuation exports a list, and every surviving false-valued
major is an empty list or one of those harmless empty iterable aliases.

Consider the possible change of goal.  If the matcher does not fire,
the goal stays numeric.  If it fires, its target test requires
\texttt{i(t,u) == a}.  The second root-opening fact therefore implies
that \texttt{t} is numeric.  By (Root scalar), \texttt{i(t,s)} and
\texttt{i(t,u)} are the same root result.  Both the Boolean-expression
and the explicit-conditional implementations consequently retain a
numeric goal.  The Boolean implementation may retain the original raw
zero alias when the transported result is false-valued; that makes no
difference to this assertion.

If the major is truthy, the generated hypothesis is the list
\texttt{m * 3}.  If the major is false-valued and \texttt{u} is truthy,
the computation expression constructs a two-element list.  A successful
opening of either hypothesis is a list.  The final recursive call then
has a numeric root-opened goal and an entirely list-valued context.  It
is a proper subtrace of the selected invocation, contradicting its
minimality.

The only remaining case is introduction.  Its expression
\texttt{i(*a,t,2)} must unpack the numeric raw goal as an iterable,
which raises a type error.  That case cannot belong to a successful
trace.  All possibilities have been excluded.\hfill$\square$

The minimization ranges over all starting floors and stacks.  It is
therefore unaffected by a recursively checked major decreasing the
global floor before its parent resumes.  The argument uses finite
completed traces, not a presumption that recursive checking always
terminates.

\subsection{Every reached command has a list-valued context}

Initially the only assumption is the list \texttt{[]}.  Suppose a
command is reached with list-valued assumptions and its major check
returns.  The lemma excludes a numeric major.  A truthy surviving major
exports a list; a computation constructs a list.  In introduction the
assignment is
\begin{verbatim}
    h = a = i(*a, t, 2)
\end{verbatim}
in mathematical order.  If its result is numeric, the later
\texttt{len(a)} raises a type error before the continuation context is
constructed, since this branch has false-valued \texttt{m}.  Such a
scalar matrix does \emph{not} insert a scalar-zero assumption.  If the
assignment yields a list, successful opening again yields a list.

Both recursive premises thus receive list-valued contexts whenever
execution reaches them.  The expression \texttt{g + [new]} constructs a
new list; a premise cannot mutate its caller's context.  Induction on
the reached calls proves the invariant.  The tuple-based explanatory
variant has the corresponding immutability property.

It follows that no numeric zero, Boolean \texttt{False}, floating zero,
truthy numeric scalar, or newly allocated numeric name can be a
successfully verified major at a reached command.  In particular a
closed scalar-zero goal is rejected by this interface.  The canonical
empty-array truth is accepted instead; mathematical completeness is
claimed for canonical encodings, not for every redundant truth alias.

\subsection{The unguarded matcher is safe on every reached continuation}

If \texttt{s} is false-valued, both matcher alternatives are nonempty
lists, so no false-valued major can match.  If \texttt{s} is truthy, the
second alternative is \texttt{False}.  A false-valued major reaching
this point is an empty list, or in the generic profile an empty string
or dictionary.  None equals \texttt{False}; a numeric major has already
been excluded.  Therefore, at every point reached after a successful
major check, the matcher has exactly the same truth value as
\begin{verbatim}
    m and m in ([s, u], not s and [t] * 3)
\end{verbatim}
even though this identity is false on unrestricted Python tuples of
values.  No sentinel and no extra source guard are needed.

The proof depends on the initial context being \texttt{[[]]}, not
\texttt{[0]}, and on premise contexts being freshly constructed.
Silently changing either detail would invalidate this argument.

\subsection{Why the relocated induction mode must be the initializer field}

The corrected selector compares arity parity with \texttt{u}.  In an
introduction branch this field is necessarily false-valued.  No
false-valued JSON value compares equal to one.  Thus a meaningful
compound matrix, of odd arity one or three, cannot receive a progress
assumption.  A literal equation has even arity two and receives
progress precisely for a zero-like mode such as \texttt{0}; an empty
array mode chooses the zero instance.  Surplus arities retain their
inherited constant interpretations.  Scalar matrices cannot survive
the shape check discussed above.

Comparing instead with the unused \texttt{s} field would be wrong.
That field is unrestricted in introduction, so choosing
\texttt{s=1} would enable progress for odd-arity compound matrices.
The distinction is semantic, not merely a preferred serialization.

In the computation branch \texttt{u} is truthy.  Its new comparison can
choose a fresh-name opening for some noncanonical scalar initializer
aliases.  This is handled by the inherited zero-instance discharge
lemma, not by claiming that every generated open hypothesis is
uniformly true.  If the generated hypothesis is $H(\vec b)$, its zero
completion $H(\vec0)$ is a computation theorem, while the old-floor
validation makes the finite family of newly admitted names $\vec b$
fresh for the retained conclusion and context.  Substituting zero for
those names and cutting in the computation theorem discharges the
extra hypothesis.  Negative fractional names in the newly admitted
interval are included separately in this finite substitution.

\subsection{Conservativity and canonical completeness}

The initial empty-array formula is already a theorem.  For example the
permitted nonempty-seed computation gives
\[
 C=0,\qquad C:=I(S(S0),S0;v.0).
\]
Use this equation as both the equality major and the premise of
transport through $F(w)\equiv(w=0)$ to obtain $0=0$.  Adding that fixed
theorem to a proof context is conservative by substitution of its
finite derivation for hypothesis uses.  No new object-level truth
predicate or logical axiom has been introduced.

For canonical certificates, substantive major commands retain their
old fields.  A computation uses the absent-major field \texttt{[]} and
its existing nonempty initializer.  An introduction uses absent major
\texttt{[]}, with \texttt{u=0} for equation progress and
\texttt{u=[]} for cases; its witness remains in \texttt{t}.
The added empty major closes without a command or name allocation, so
the remaining canonical proof simulation is unchanged.  First applying
the existing certified-zero repair supplies every unrestricted
computation instance.  This gives a finite certificate for every
canonical derivation of the inherited native calculus.

Conversely, the invariant makes each reached matcher a legitimate
guarded matcher.  The remaining commands use the inherited
transport, closure, paired-introduction/elimination, and zero-instance
computation arguments, with the corrected equation-only mode test.
Discharging the initial $0=0$ theorem yields a derivation in that same
calculus.  Its established whole-domain iPRA interpretations therefore
transfer without changing the numerical domain, equality, or object
language.  The new result concerns a proof-checker presentation, not a
new primitive-recursive arithmetic foundation.

\subsection{Raw JSON types and the generic profile}

The strict profiles recognize only Python lists as containers.  The
short predicate \texttt{a * 0 == []} has this effect for JSON values:
lists pass; strings and numbers give different zero products;
dictionaries and null raise.  Its numerical fallback rejects strings.
The root invariant covers all numbers and Booleans that reach it,
including nonfinite values accepted by the default Python loader.

The generic profile has the inherited additional aliases.  An empty
string and an empty dictionary open to \texttt{[]}; a dictionary with
only the empty-string key opens to \texttt{[[]]}, with its value ignored.
A truthy dictionary cannot survive export as a major because
\texttt{m * 3} raises.  Nonempty strings and nonempty dictionary keys
have no finite successful opening.  Actual CPython rejects them through
its recursion limit and subsequent type failure, while an ideal
unlimited-stack execution can diverge.  The raw soundness argument
above covers every finite successful run of this profile; it does not
assert ideal termination on every malformed finite JSON input.

Larger arrays remain arrays throughout opening, irrespective of their
ignored annotations or depth increments.  They therefore do not
threaten the scalar-versus-list invariant.  Their formula meanings and
freshness behavior remain subject to the inherited sorted-decoder
proof.  Physical time, memory, and stack limits can reject sufficiently
large genuine proofs in every profile; the completeness claim has the
same adequate-resource, canonical-proof interpretation as its
predecessors.


\clearpage
\begin{thebibliography}{99}
\bibitem{corpus}
Vladimir Reshetnikov, \emph{ProveIt, docs/drafts}, commit
\file{b9c4d74a0018bdafc5bc06e755a0145f4ed2b8d6}, inspected 10 November 2026.
\url{https://github.com/VladimirReshetnikov/ProveIt/tree/b9c4d74a0018bdafc5bc06e755a0145f4ed2b8d6/docs/drafts}.
Archive paths, byte lengths and content hashes are recorded in the supplied
\file{evidence/corpus_manifest.json}.
\bibitem{neutral}
\emph{Equality from Monus and Neutral Major Premises},
\file{ipra_monus_and_neutral_major_2026-11-10.zip}, in \cite{corpus}.
The 542-byte incumbent and the neutral-major proof are reproduced here.
\bibitem{uniform}
\emph{Uniform Zero and Constructed Premises for iPRA},
\file{ipra_uniform_zero_2026-11-10.zip}, in \cite{corpus}.
Source of the matched lexical-floor affine checker and local-scope producer.
\bibitem{selfseed}
\emph{Scalar Self-Seeded Arithmetic and the iPRA Checker Frontier},
\file{ipra_selfseed_scalar_2026-11-10.zip}, in \cite{corpus}.
Its self-seeded completion explicitly assumes successor injectivity.
\bibitem{certified}
\emph{A Certified Zero Removes a Computation Opcode},
\file{ipra_certified_zero_2026-11-10.zip}, in \cite{corpus}.
Source of the certified-zero normalization and the attributed native and
arithmetic appendices reproduced here.
\bibitem{arity}
\emph{An Arity Calculus for iPRA},
\file{ipra_arity_calculus_2026-11-10.zip}, in \cite{corpus}.
The historical arity proof is preserved in \cite{certified},
\file{reference/inherited_arity_checker.tex}.
\bibitem{constructed}
\file{ipra_constructed_computation_2026-11-10.zip}, in \cite{corpus}.
The constructed-computation antecedent used by the arity representation.
\bibitem{generated}
Generated-assumption antecedents, preserved and attributed in the
\file{reference/} material of \cite{certified,constructed}.
The present report relies on the reproduced statements and proofs,
not on an unavailable external proof checker.
\bibitem{transient}
\emph{A One-State Recursor for iPRA},
\file{ipra_transient_seed_2026-11-10.zip}, in \cite{corpus}.
\bibitem{erasure}
Scalar predecessor development, reproduced in \cite{certified} as
\file{reference/scalar_predecessor.tex}. The exact reproduced source
hash is retained in this package's copied-source inventory.
\bibitem{forward}
Literal-equation induction and equality-reflection development,
reproduced in \cite{certified} as \file{reference/equality_reflection.tex}.
\bibitem{progress}
Pairing and primitive-recursion compilation, reproduced in
\cite{certified} as \file{reference/primitive_recursion.tex}.
\bibitem{allstarts}
\emph{Induction from Arbitrary Starting Points}, preserved in
\file{ipra_research_2026-11-10.zip},
\file{reference/audit/prior_math/all_starts.tex}, in \cite{corpus}.
It already contains deterministic progression/closure introduction, the
separate Step/Base predecessor calculus, and its all-starts fusion.
\end{thebibliography}
\end{document}
